% =====================================================================
%  GOLDEN REASONING -- MASTER DOCUMENT
%  Sections 1-2 assembled 2026-08-01.
%
%  SOURCES.  This file is the concatenation of two standalone section
%  files, each of which carries its own full provenance and correction
%  ledger in its header:
%      section-soft-horizons.tex   -> Section 1
%      section-two-channel.tex     -> Section 2 (revision 2)
%  Those ledgers are NOT reproduced here.  Keep the standalone files:
%  they are the authoritative record of who originated what, of the
%  correction log, and of citation-verification status.  Consult them
%  before posting.
%
%  BIBLIOGRAPHY.  The two sections shared four references (AS15, GJW99,
%  VMS79, VY85); the merged list is their union, 35 entries, deduped.
%  The file is SELF-SWITCHING: if biblatex is installed it is used,
%  with the .bib written at compile time by filecontents; otherwise a
%  manual thebibliography is used instead.  Both are generated from the
%  same 35-key list and their key sets were checked equal.  To force
%  one path, edit the \IfFileExists test below.
%
%  NOTE ON VERIFICATION.  The fallback (thebibliography) path is
%  compile-verified.  The biblatex path is NOT -- biblatex and biber
%  were unavailable in the environment where this file was assembled.
%  If biber errors, the usual causes are (a) needing a biber run
%  between pdflatex runs, or (b) backend=biber vs bibtex mismatch.
%
%  CROSS-REFERENCES.  Section 2's references to "Section One" have been
%  converted to \ref{sh:sec} in this merged file.  Label prefixes are
%  disjoint: sh: for Section 1, tc: for Section 2.
% =====================================================================

\documentclass[11pt]{article}

\usepackage[margin=1.1in]{geometry}
\usepackage{amsmath,amssymb,amsthm,mathtools}
\usepackage{booktabs}
\usepackage{array}
\usepackage{enumitem}
\usepackage{xcolor}
\definecolor{grlink}{rgb}{0.10,0.20,0.55}
\usepackage[colorlinks=true,linkcolor=grlink,citecolor=grlink,urlcolor=grlink]{hyperref}

% ---- bibliography database (written at compile time) ----------------
\begin{filecontents*}[overwrite]{golden-reasoning.bib}
@article{AD13, author={Aurzada, Frank and Dereich, Steffen}, title={Universality of the asymptotics of the one-sided exit problem for integrated processes}, journal={Ann. Inst. Henri Poincar\'e Probab. Stat.}, volume={49}, year={2013}, pages={236--251}}
@incollection{AS15, author={Aurzada, Frank and Simon, Thomas}, title={Persistence probabilities and exponents}, booktitle={L\'evy Matters V}, series={Lecture Notes in Mathematics}, volume={2149}, publisher={Springer}, address={Cham}, year={2015}, pages={183--224}}
@article{BBM05, author={Ben Arous, G\'erard and Bogachev, Leonid V. and Molchanov, Stanislav A.}, title={Limit theorems for sums of random exponentials}, journal={Probab. Theory Relat. Fields}, volume={132}, number={4}, year={2005}, pages={579--612}}
@article{BMS13, author={Bray, Alan J. and Majumdar, Satya N. and Schehr, Gr\'egory}, title={Persistence and first-passage properties in nonequilibrium systems}, journal={Adv. Phys.}, volume={62}, number={3}, year={2013}, pages={225--361}}
@article{Bha82, author={Bhattacharya, Rabi N.}, title={On the functional central limit theorem and the law of the iterated logarithm for Markov processes}, journal={Z. Wahrscheinlichkeitstheorie verw. Gebiete}, volume={60}, year={1982}, pages={185--201}}
@book{Bov06, author={Bovier, Anton}, title={Statistical Mechanics of Disordered Systems: A Mathematical Perspective}, publisher={Cambridge University Press}, year={2006}}
@book{CL67, author={Cram\'er, Harald and Leadbetter, M. Ross}, title={Stationary and Related Stochastic Processes}, publisher={Wiley}, address={New York}, year={1967}}
@article{DW15, author={Denisov, Denis and Wachtel, Vitali}, title={Exit times for integrated random walks}, journal={Ann. Inst. Henri Poincar\'e Probab. Stat.}, volume={51}, number={1}, year={2015}, pages={167--193}}
@article{Der80, author={Derrida, Bernard}, title={Random-energy model: limit of a family of disordered models}, journal={Phys. Rev. Lett.}, volume={45}, year={1980}, pages={79--82}}
@article{Der81, author={Derrida, Bernard}, title={Random-energy model: an exactly solvable model of disordered systems}, journal={Phys. Rev. B}, volume={24}, number={5}, year={1981}, pages={2613--2626}}
@book{Dur19, author={Durrett, Rick}, title={Probability: Theory and Examples}, edition={5}, publisher={Cambridge University Press}, year={2019}}
@book{FW12, author={Freidlin, Mark I. and Wentzell, Alexander D.}, title={Random Perturbations of Dynamical Systems}, edition={3}, publisher={Springer}, year={2012}}
@article{GJW99, author={Groeneboom, Piet and Jongbloed, Geurt and Wellner, Jon A.}, title={Integrated Brownian motion, conditioned to be positive}, journal={Ann. Probab.}, volume={27}, number={3}, year={1999}, pages={1283--1303}}
@article{Gol71, author={Goldman, Malcolm}, title={On the first passage of the integrated Wiener process}, journal={Ann. Math. Statist.}, volume={42}, year={1971}, pages={2150--2155}}
@article{Gom25, author={Gompertz, Benjamin}, title={On the nature of the function expressive of the law of human mortality, and on a new mode of determining the value of life contingencies}, journal={Philos. Trans. R. Soc. Lond.}, volume={115}, year={1825}, pages={513--583}}
@article{Hor67, author={H\"ormander, Lars}, title={Hypoelliptic second order differential equations}, journal={Acta Math.}, volume={119}, year={1967}, pages={147--171}}
@article{IW94, author={Isozaki, Yasuki and Watanabe, Shinzo}, title={An asymptotic formula for the Kolmogorov diffusion and a refinement of Sinai's estimates for the integral of Brownian motion}, journal={Proc. Japan Acad. Ser. A}, volume={70}, number={9}, year={1994}, pages={271--276}}
@book{KT81, author={Karlin, Samuel and Taylor, Howard M.}, title={A Second Course in Stochastic Processes}, publisher={Academic Press}, year={1981}}
@article{Kol34, author={Kolmogorov, Andrei N.}, title={Zuf\"allige Bewegungen (zur Theorie der Brownschen Bewegung)}, journal={Ann. of Math. (2)}, volume={35}, year={1934}, pages={116--117}}
@article{Lac91, author={Lachal, Aim\'e}, title={Sur le premier instant de passage de l'int\'egrale du mouvement brownien}, journal={Ann. Inst. H. Poincar\'e Probab. Statist.}, volume={27}, year={1991}, pages={385--405}}
@book{Lan04, author={Lando, David}, title={Credit Risk Modeling: Theory and Applications}, publisher={Princeton University Press}, year={2004}}
@article{MK04, author={Molchan, George and Khokhlov, Andrei}, title={Small values of the maximum for the integral of fractional Brownian motion}, journal={J. Stat. Phys.}, volume={114}, year={2004}, pages={923--945}}
@article{MP01, author={Milevsky, Moshe A. and Promislow, S. David}, title={Mortality derivatives and the option to annuitise}, journal={Insurance Math. Econom.}, volume={29}, number={3}, year={2001}, pages={299--318}}
@article{MV12, author={M\'el\'eard, Sylvie and Villemonais, Denis}, title={Quasi-stationary distributions and population processes}, journal={Probab. Surv.}, volume={9}, year={2012}, pages={340--410}}
@article{Mar49, author={Maruyama, Gisiro}, title={The harmonic analysis of stationary stochastic processes}, journal={Mem. Fac. Sci. Ky\=ush\=u Univ. Ser. A}, volume={4}, year={1949}, pages={45--106}}
@article{McK63, author={McKean, Jr., Henry P.}, title={A winding problem for a resonator driven by a white noise}, journal={J. Math. Kyoto Univ.}, volume={2}, year={1963}, pages={227--235}}
@article{Mol99, author={Molchan, George M.}, title={Maximum of a fractional Brownian motion: probabilities of small values}, journal={Comm. Math. Phys.}, volume={205}, year={1999}, pages={97--111}}
@article{RVY06, author={Roynette, Bernard and Vallois, Pierre and Yor, Marc}, title={Some penalisations of the Wiener measure}, journal={Jpn. J. Math.}, volume={1}, year={2006}, pages={263--290}}
@book{RY99, author={Revuz, Daniel and Yor, Marc}, title={Continuous Martingales and Brownian Motion}, edition={3}, publisher={Springer}, year={1999}}
@article{Ric44, author={Rice, Stephen O.}, title={Mathematical analysis of random noise}, journal={Bell System Tech. J.}, volume={23}, year={1944}, pages={282--332}}
@article{Ric45, author={Rice, Stephen O.}, title={Mathematical analysis of random noise (concluded)}, journal={Bell System Tech. J.}, volume={24}, year={1945}, pages={46--156}}
@article{Sin92, author={Sinai, Yakov G.}, title={Distribution of some functionals of the integral of a random walk}, journal={Theoret. and Math. Phys.}, volume={90}, year={1992}, pages={219--241}, note={translated from Teoret. Mat. Fiz. \textbf{90} (1992), no. 3, 323--353}}
@article{VMS79, author={Vaupel, James W. and Manton, Kenneth G. and Stallard, Eric}, title={The impact of heterogeneity in individual frailty on the dynamics of mortality}, journal={Demography}, volume={16}, year={1979}, pages={439--454}}
@article{VY85, author={Vaupel, James W. and Yashin, Anatoli I.}, title={Heterogeneity's ruses: some surprising effects of selection on population dynamics}, journal={Amer. Statist.}, volume={39}, year={1985}, pages={176--185}}
@article{Wah78, author={Wahba, Grace}, title={Improper priors, spline smoothing and the problem of guarding against model errors in regression}, journal={J. Roy. Statist. Soc. Ser. B}, volume={40}, year={1978}, pages={364--372}}
\end{filecontents*}

\begin{filecontents*}[overwrite]{gr-manualbib.tex}
\begin{thebibliography}{99}\small
\bibitem[AD13]{AD13} F.~Aurzada and S.~Dereich,
\emph{Universality of the asymptotics of the one-sided exit problem for
integrated processes}, Ann.\ Inst.\ Henri Poincar\'e Probab.\ Stat.\
\textbf{49} (2013), 236--251.

\bibitem[AS15]{AS15} F.~Aurzada and T.~Simon,
\emph{Persistence probabilities and exponents}, in: L\'evy Matters V,
Lecture Notes in Mathematics 2149, Springer, Cham, 2015, pp.~183--224.

\bibitem[BBM05]{BBM05} G.~Ben~Arous, L.~V.~Bogachev, and
S.~A.~Molchanov, \emph{Limit theorems for sums of random exponentials},
Probab.\ Theory Relat.\ Fields \textbf{132} (2005), no.~4, 579--612.

\bibitem[Bha82]{Bha82} R.~N.~Bhattacharya,
\emph{On the functional central limit theorem and the law of the iterated
logarithm for Markov processes},
Z.~Wahrscheinlichkeitstheorie verw.\ Gebiete \textbf{60} (1982), 185--201.

\bibitem[BMS13]{BMS13} A.~J.~Bray, S.~N.~Majumdar, and G.~Schehr,
\emph{Persistence and first-passage properties in nonequilibrium
systems}, Adv.\ Phys.\ \textbf{62} (2013), 225--361.

\bibitem[Bov06]{Bov06} A.~Bovier,
\emph{Statistical Mechanics of Disordered Systems: A Mathematical
Perspective}, Cambridge University Press, 2006.

\bibitem[CL67]{CL67} H.~Cram\'er and M.~R.~Leadbetter,
\emph{Stationary and Related Stochastic Processes}, Wiley, New York,
1967.

\bibitem[Der80]{Der80} B.~Derrida,
\emph{Random-energy model: limit of a family of disordered models},
Phys.\ Rev.\ Lett.\ \textbf{45} (1980), 79--82.

\bibitem[Der81]{Der81} B.~Derrida,
\emph{Random-energy model: an exactly solvable model of disordered
systems}, Phys.\ Rev.\ B \textbf{24} (1981), no.~5, 2613--2626.

\bibitem[Dur19]{Dur19} R.~Durrett,
\emph{Probability: Theory and Examples}, 5th ed., Cambridge University
Press, 2019.

\bibitem[DW15]{DW15} D.~Denisov and V.~Wachtel,
\emph{Exit times for integrated random walks}, Ann.\ Inst.\ Henri
Poincar\'e Probab.\ Stat.\ \textbf{51} (2015), 167--193.

\bibitem[FW12]{FW12} M.~I.~Freidlin and A.~D.~Wentzell,
\emph{Random Perturbations of Dynamical Systems}, 3rd ed., Springer, 2012.

\bibitem[GJW99]{GJW99} P.~Groeneboom, G.~Jongbloed, and J.~A.~Wellner,
\emph{Integrated Brownian motion, conditioned to be positive},
Ann.\ Probab.\ \textbf{27} (1999), 1283--1303.

\bibitem[Gol71]{Gol71} M.~Goldman,
\emph{On the first passage of the integrated Wiener process},
Ann.\ Math.\ Statist.\ \textbf{42} (1971), 2150--2155.

\bibitem[Gom25]{Gom25} B.~Gompertz,
\emph{On the nature of the function expressive of the law of human
mortality, and on a new mode of determining the value of life
contingencies}, Philos.\ Trans.\ R.\ Soc.\ Lond.\ \textbf{115} (1825),
513--583.

\bibitem[H\"or67]{Hor67} L.~H\"ormander,
\emph{Hypoelliptic second order differential equations},
Acta Math.\ \textbf{119} (1967), 147--171.

\bibitem[IW94]{IW94} Y.~Isozaki and S.~Watanabe,
\emph{An asymptotic formula for the Kolmogorov diffusion and a refinement
of Sinai's estimates for the integral of Brownian motion},
Proc.\ Japan Acad.\ Ser.\ A \textbf{70} (1994), 271--276.

\bibitem[Kol34]{Kol34} A.~N.~Kolmogorov,
\emph{Zuf\"allige Bewegungen (zur Theorie der Brownschen Bewegung)},
Ann.\ of Math.\ (2) \textbf{35} (1934), 116--117.

\bibitem[KT81]{KT81} S.~Karlin and H.~M.~Taylor,
\emph{A Second Course in Stochastic Processes}, Academic Press, 1981.

\bibitem[Lac91]{Lac91} A.~Lachal,
\emph{Sur le premier instant de passage de l'int\'egrale du mouvement
brownien}, Ann.\ Inst.\ H.\ Poincar\'e Probab.\ Statist.\ \textbf{27}
(1991), 385--405.

\bibitem[Lan04]{Lan04} D.~Lando,
\emph{Credit Risk Modeling: Theory and Applications}, Princeton
University Press, 2004.

\bibitem[Mar49]{Mar49} G.~Maruyama,
\emph{The harmonic analysis of stationary stochastic processes},
Mem.\ Fac.\ Sci.\ Ky\=ush\=u Univ.\ Ser.\ A \textbf{4} (1949), 45--106.

\bibitem[McK63]{McK63} H.~P.~McKean, Jr.,
\emph{A winding problem for a resonator driven by a white noise},
J.\ Math.\ Kyoto Univ.\ \textbf{2} (1963), 227--235.

\bibitem[MK04]{MK04} G.~Molchan and A.~Khokhlov,
\emph{Small values of the maximum for the integral of fractional
Brownian motion}, J.\ Stat.\ Phys.\ \textbf{114} (2004), 923--945.

\bibitem[Mol99]{Mol99} G.~M.~Molchan,
\emph{Maximum of a fractional Brownian motion: probabilities of small
values}, Comm.\ Math.\ Phys.\ \textbf{205} (1999), 97--111.

\bibitem[MP01]{MP01} M.~A.~Milevsky and S.~D.~Promislow,
\emph{Mortality derivatives and the option to annuitise},
Insurance Math.\ Econom.\ \textbf{29} (2001), 299--318.

\bibitem[MV12]{MV12} S.~M\'el\'eard and D.~Villemonais,
\emph{Quasi-stationary distributions and population processes},
Probab.\ Surv.\ \textbf{9} (2012), 340--410.

\bibitem[Ric44]{Ric44} S.~O.~Rice,
\emph{Mathematical analysis of random noise}, Bell System Tech.\ J.\
\textbf{23} (1944), 282--332.

\bibitem[Ric45]{Ric45} S.~O.~Rice,
\emph{Mathematical analysis of random noise (concluded)}, Bell System
Tech.\ J.\ \textbf{24} (1945), 46--156.

\bibitem[RVY06]{RVY06} B.~Roynette, P.~Vallois, and M.~Yor,
\emph{Some penalisations of the Wiener measure}, Jpn.\ J.\ Math.\
\textbf{1} (2006), 263--290.

\bibitem[RY99]{RY99} D.~Revuz and M.~Yor,
\emph{Continuous Martingales and Brownian Motion}, 3rd ed., Springer,
1999.

\bibitem[Sin92]{Sin92} Ya.~G.~Sinai,
\emph{Distribution of some functionals of the integral of a random walk},
Theoret.\ and Math.\ Phys.\ \textbf{90} (1992), 219--241; translated
from Teoret.\ Mat.\ Fiz.\ \textbf{90} (1992), no.~3, 323--353.

\bibitem[VMS79]{VMS79} J.~W.~Vaupel, K.~G.~Manton, and E.~Stallard,
\emph{The impact of heterogeneity in individual frailty on the dynamics
of mortality}, Demography \textbf{16} (1979), 439--454.

\bibitem[VY85]{VY85} J.~W.~Vaupel and A.~I.~Yashin,
\emph{Heterogeneity's ruses: some surprising effects of selection on
population dynamics}, Amer.\ Statist.\ \textbf{39} (1985), 176--185.

\bibitem[Wah78]{Wah78} G.~Wahba,
\emph{Improper priors, spline smoothing and the problem of guarding
against model errors in regression},
J.\ Roy.\ Statist.\ Soc.\ Ser.\ B \textbf{40} (1978), 364--372.

\end{thebibliography}
\end{filecontents*}

\IfFileExists{biblatex.sty}{%
  \usepackage[backend=biber,style=alphabetic,maxbibnames=99,giveninits=true]{biblatex}%
  \addbibresource{golden-reasoning.bib}%
  \newcommand{\PrintRefs}{\printbibliography[title={References}]}%
}{%
  \newcommand{\PrintRefs}{\input{gr-manualbib.tex}}%
}

\numberwithin{equation}{section}

\theoremstyle{plain}
\newtheorem{theorem}{Theorem}[section]
\newtheorem{proposition}[theorem]{Proposition}
\newtheorem{lemma}[theorem]{Lemma}
\newtheorem{corollary}[theorem]{Corollary}
\newtheorem{claim}[theorem]{Claim}

\theoremstyle{definition}
\newtheorem{definition}[theorem]{Definition}
\newtheorem{ansatz}[theorem]{Ansatz}
\newtheorem{convention}[theorem]{Convention}
\newtheorem{constraint}{Constraint}

\theoremstyle{remark}
\newtheorem{remark}[theorem]{Remark}
\newtheorem{example}[theorem]{Example}

\newcommand{\E}{\mathbb{E}}
\newcommand{\PP}{\mathbb{P}}
\newcommand{\R}{\mathbb{R}}
\DeclareMathOperator{\Var}{Var}
\DeclareMathOperator{\Cov}{Cov}
\DeclareMathOperator{\Leb}{Leb}
\newcommand{\eqd}{\overset{d}{=}}
\newcommand{\one}{\mathbf{1}}
\newcommand{\Imm}{\mathcal{I}}
\newcommand{\Fc}{\mathcal{F}}
\newcommand{\Lc}{\mathcal{L}}
\newcommand{\hm}{\mathfrak{h}}
\newcommand{\Ac}{\mathcal{A}}
\newcommand{\Sc}{\mathsf{S}}

\allowdisplaybreaks

\title{Golden Reasoning\\[4pt]
\large Soft Horizons, Survivorship, and the Two-Channel Muffled Hazard}
\author{}
\date{August 2026}

\begin{document}
\maketitle
\tableofcontents
\bigskip

\section{Soft horizons and survivorship in stochastic hazard dynamics}\label{sh:sec}

\begin{quote}\small
This section is organized as the buildup to a single question. Let $T$ be
the death time of a survival system whose hazard rate $h_t$ has negative
logarithmic derivative $\mu_t=-\frac{d}{dt}\log h_t$ driven by Brownian
motion, and consider the \emph{survivorship interpolation}: the family of
conditioned laws $\PP^{(t)}=\PP(\,\cdot\mid T>t)$ joining the unbiased law
$\PP^{(0)}=\PP$ --- the ensemble of systems that have survived at least
time zero --- to the ensembles of ever-longer survivors. What is the
structure of this interpolation, and what is its asymptotic behavior as
$t\to\infty$? The state-space theory comes first because it is what the
answer displaces: the state process is a Kolmogorov diffusion whose
initial data form a gauge orbit invisible to asymptotics
(\S\ref{sh:sub:model}); immortality is impossible for every initial
condition, yet life expectancy is infinite, the survival tail carrying the
persistence exponent $\tfrac14$ of integrated Brownian motion
(\S\ref{sh:sub:verdict}); and genuine horizons in the initial condition
exist exactly when the driver's scale function is bounded, the horizon
profile being the normalized scale function itself --- degenerate for the
Brownian driver, an explicit Gaussian sigmoid for the unstable
Ornstein--Uhlenbeck driver (\S\ref{sh:sub:horizon}). The verdict
(\S\ref{sh:sub:interp}) is that the horizon missing from state space is
real but lives in the ensemble. The interpolation is a monotone flow of
upward tilts converging, at polynomial rate $s/4t$ on each window of
length $s$, to the pure Doob transform by the persistence-harmonic
function: a limit locally absolutely continuous on survivors yet globally
singular, whose generator acquires the universal survivorship drift
$\sigma^2\partial_\mu\log\hm$, reducing at zero credit to
$\sigma^2/2\mu$, under which the credit passes from the recurrent to the
transient class: the P\'olya $2\mathrm{D}\to3\mathrm{D}$ promotion,
effected by conditioning rather than by dimension. Across driver classes
the interpolation has
exactly three asymptotic modes --- saturation, transformation,
concentration --- in bijection with the horizon trichotomy, the driftless
model realizing the critical mode.
\end{quote}

\subsection*{Overview and status of results}

Throughout, statements are tagged by their epistemic status:
\textbf{[P]}~proved here in full; \textbf{[C]}~classical or cited;
\textbf{[A]}~conditional on the soft-wall reduction, Ansatz~\ref{sh:ansatz};
\textbf{[S]}~derived at the level of a large-deviation or scaling sketch;
\textbf{[O]}~open. Inside statement headers the brackets are dropped, as
in \emph{(No viability; \textbf{P})}, the surrounding optional argument
supplying them. The load-bearing structure is:
the model is exactly a Kolmogorov diffusion with multiplicative-exponential
killing \textbf{[P]}; killed at exponent level, it is the persistence problem
for integrated Brownian motion \textbf{[A]}, whose exponent $\tfrac14$ is a
theorem \textbf{[C]}; certain death for the driftless model is proved
unconditionally \textbf{[P]}; the trichotomy governing when genuine horizons
exist is proved \textbf{[P]}, its critical recurrent case proved under mild
hypotheses \textbf{[P]}; the solvable gradual horizon is exact \textbf{[P]};
the sharp survival rate under adverse drift is identified variationally
\textbf{[S]} with rigorous one-sided bounds \textbf{[P]}. For the
destination subsection: monotonicity of the survivorship interpolation is
proved \textbf{[P]}; its critical limit is identified as the pure Doob
transform \textbf{[C]}, transferred to the soft-kill model \textbf{[A]};
the survivorship drift, its positivity, and its wall asymptotics are
derived \textbf{[P]} from cited inputs; local absolute
continuity together with global singularity is established
\textbf{[P/C]}; and
the three interpolation modes are established at levels
\textbf{[P]}/\textbf{[C]}/\textbf{[S]} respectively. A reader wanting the
results at a glance, rather than in order of proof, may turn directly to
the reference sheet of \S\ref{sh:sub:sheet}.

\subsection*{Relation to the literature}

Four bodies of work border this section, and the reader deserves a
precise map of what is classical, what is adjacent, and what is claimed
as new.

\emph{The classical layer.} The state process is Kolmogorov's 1934
diffusion \cite{Kol34}, the primordial hypoelliptic example
\cite{Hor67}. Its persistence theory is a mature subject with a clean
lineage --- McKean \cite{McK63}, Goldman \cite{Gol71}, Lachal
\cite{Lac91}, Sinai \cite{Sin92}, Isozaki--Watanabe \cite{IW94},
Denisov--Wachtel \cite{DW15} --- surveyed in \cite{AS15} and, as the
\emph{random acceleration process}, possessing an extensive physics
literature reviewed in \cite{BMS13}. The conditioned (never-hitting)
process, our limit object in \S\ref{sh:sub:interp}, is constructed by
Groeneboom, Jongbloed and Wellner \cite{GJW99}. None of this is ours;
we import it.

\emph{The adjacent frameworks.} Our survivorship interpolation belongs
to two established genres. First, Theorem~\ref{sh:thm:Q} is a
\emph{penalisation} statement in the sense of Roynette--Vallois--Yor
\cite{RVY06}: a limit of Wiener-type measures reweighted by a
time-dependent functional, identified as a martingale change of measure
--- here the weight is survival itself, and the polynomial tail places
it in the clockless class. Second, the mode classification of
Theorem~\ref{sh:thm:modes} sits against the theory of quasi-stationary
distributions and $Q$-processes surveyed in \cite{MV12}: classical QSD
theory is at home in the exponential-tail class (our modes (I) and
(III)), while the critical polynomial mode admits no proper QSD in fixed
coordinates, its Yaglom object being self-similar. We use both
frameworks as positioning, not as inputs; the proofs here are
self-contained given the cited persistence theory.

\emph{The applied layer: what this is a model of.} Stochastic hazard
rates are not new, and the wrapper $S_t=\exp(-\int_0^th)$ with $h$ a
diffusion is entirely standard: it is the reduced-form, or
intensity-based, framework of credit risk \cite{Lan04} and of stochastic
mortality, where the force of mortality is routinely modelled by a
diffusion, the mean-reverting Brownian Gompertz specification of Milevsky
and Promislow \cite{MP01} being a canonical instance. What distinguishes
the present model is not the wrapper but \emph{where the noise is
placed}: standard practice diffuses $h$ or $\log h$, whereas we diffuse
$-\tfrac{d}{dt}\log h$, one integration further out. The consequences are
structural rather than cosmetic --- $\log h$ becomes $C^1$, the state
becomes degenerate and $3/2$-self-similar, and the governing persistence
exponent moves from $\tfrac12$ to $\tfrac14$ (Remark~\ref{sh:rem:iwp},
\S\ref{sh:subsub:tail}). We stress the point because it locates the
novelty precisely, and disclaims the rest.

\emph{The conceptual ancestor.} The thesis of \S\ref{sh:sub:interp} ---
that selection on survival makes the observed ensemble diverge
systematically from the unbiased one --- is not new in substance. It is
the frailty phenomenon of Vaupel, Manton and Stallard \cite{VMS79} and
Vaupel and Yashin \cite{VY85}: when the frail are removed first the
population hazard falls below the individual hazard, the gap widening
with age, producing mortality deceleration at the population level with
no corresponding deceleration for any individual. We claim no priority
over that insight, and regard this section as exhibiting its
\emph{critical dynamical} analogue. The differences are structural and
are the point. Frailty is a static random variable fixed at birth and the
analysis is a mixture computation over a population; here the randomness
is a driver evolving in time, and the object is a limit of path measures
singular with respect to the unbiased law
(Theorem~\ref{sh:thm:singular}). Where the frailty literature computes
how an observed hazard curve bends, the interpolation of
\S\ref{sh:sub:interp} identifies the limiting law of the survivor's
entire trajectory, together with the drift \eqref{sh:eq:G} that
conditioning manufactures.

\emph{What is claimed as new.} To our knowledge, the following are new
at least in assembly, and we flag them as this section's contribution:
the log-derivative placement of the noise, with its gauge analysis
(\S\ref{sh:subsub:dimensional}) and the two-sided soft-wall analysis
whose remaining gap is located exactly (Ansatz~\ref{sh:ansatz}); the
scale-function horizon trichotomy with mortal critical case
(\S\ref{sh:subsub:scale}); the adverse-drift survival rate
$3\nu^2/8\sigma^2$ with its front-load-and-coast optimal profile
(Proposition~\ref{sh:prop:advdrift}); and the survivorship
interpolation read as a monotone flow with the tail index as its
linear-response coefficient, classified across drivers into saturation,
transformation and concentration (\S\ref{sh:sub:interp}). Where a
statement merely re-expresses known results in this frame, the tags and
citations say so.

\begin{convention}\label{sh:conv}
For positive functions, $f\sim g$ means $f/g\to1$; $f\asymp g$ means
$0<\liminf f/g\le\limsup f/g<\infty$; $f\lesssim g$ means $f\le Cg$.
Hatted quantities are nondimensionalized by the intrinsic timescale of
Definition~\ref{sh:def:units}. $\Phi$ is the standard normal distribution
function; $\one_A$ is an indicator; $\Leb$ is Lebesgue measure. All processes
are defined on a filtered probability space satisfying the usual conditions,
and $W$ denotes a standard one-dimensional Brownian motion.
\end{convention}

% ---------------------------------------------------------------------
\subsection{The model and its invariances}\label{sh:sub:model}

\subsubsection{Setup: hazard, credit, and the Kolmogorov diffusion}
\label{sh:subsub:setup}

\begin{definition}[Model]\label{sh:def:model}
Fix $\sigma>0$ and initial data $(h_0,\mu_0,D_0)\in(0,\infty)\times\R\times\R$.
Define the \emph{improvement rate}, \emph{credit}, \emph{hazard}, and
\emph{cumulative hazard} by
\begin{equation}\label{sh:eq:model}
\mu_t=\mu_0+\sigma W_t,\qquad
D_t=D_0+\int_0^t\mu_s\,ds,\qquad
h_t=h_0e^{-(D_t-D_0)},\qquad
\Lambda_t=\int_0^t h_u\,du,
\end{equation}
and write $x_t=\log h_t$, so that $\mu_t=-\dot x_t$. We also set
$I_t=\int_0^t W_s\,ds$, so that $D_t=D_0+\mu_0t+\sigma I_t$.
\end{definition}

\begin{definition}[Death time; Cox construction]\label{sh:def:cox}
Let $E\sim\mathrm{Exp}(1)$ be independent of $W$ and define
$T:=\inf\{t\ge0:\Lambda_t\ge E\}\in(0,\infty]$. Then
$\PP(T>t\mid\Fc^W_\infty)=e^{-\Lambda_t}$ and
$\PP(T=\infty\mid\Fc^W_\infty)=e^{-\Lambda_\infty}$, with the convention
$e^{-\infty}=0$. The \emph{survival function} is
$S_t=\PP(T>t\mid\Fc^W_\infty)=e^{-\Lambda_t}$.
\end{definition}

\begin{definition}[Viability]\label{sh:def:viab}
The \emph{viability event} is $\Imm:=\{\Lambda_\infty<\infty\}$. On $\Imm$
the conditional survival probability $S_\infty=e^{-\Lambda_\infty}$ is
strictly positive; off $\Imm$ it vanishes. Thus
\begin{equation}\label{sh:eq:twotier}
\PP(T=\infty)\;=\;\E\!\left[e^{-\Lambda_\infty}\one_\Imm\right]
\;\le\;\PP(\Imm),
\end{equation}
and we distinguish throughout between \emph{viability} ($\PP(\Imm)$, the
probability that immortality is achievable) and \emph{immortality proper}
($\PP(T=\infty)$).
\end{definition}

The pair $(\mu,D)$ is the \emph{Kolmogorov diffusion} \cite{Kol34}:
\begin{equation}\label{sh:eq:kolm}
d\mu_t=\sigma\,dW_t,\qquad dD_t=\mu_t\,dt,\qquad
\Lc=\tfrac{\sigma^2}{2}\partial_\mu^2+\mu\,\partial_D .
\end{equation}

\begin{lemma}[Regularity]\label{sh:lem:hypo}
The free process $(\mu_t,D_t)$ is Gaussian with mean
$(\mu_0,\,D_0+\mu_0t)$ and covariance
$\bigl(\begin{smallmatrix}\sigma^2t&\sigma^2t^2/2\\
\sigma^2t^2/2&\sigma^2t^3/3\end{smallmatrix}\bigr)$;
in particular it admits the explicit smooth transition density written down
by Kolmogorov \cite{Kol34}. More generally, $\Lc$ is hypoelliptic: with
$X_1=\sigma\partial_\mu$ and $X_0=\mu\partial_D$ one has
$[X_1,X_0]=\sigma\partial_D$, so H\"ormander's bracket condition holds at
first order \cite{Hor67}. Consequently the process killed at rate $e^{x}$
retains smooth densities, which is the regularity relevant to the killed
problem of \S\ref{sh:subsub:reduction}.
\end{lemma}

\begin{proof}
Gaussianity and the moments follow from
Lemma~\ref{sh:lem:secondorder} below and linearity of
\eqref{sh:eq:model} in $W$; the bracket computation is immediate.
\end{proof}

Three structural remarks frame everything that follows.

\begin{remark}[No boundary in $h$]\label{sh:rem:noboundary}
The logarithmic parametrization makes $h>0$ automatic: there is no absorbing
or reflecting boundary in the hazard itself, and the entire problem concerns
the growth of the integrated driver.
\end{remark}

\begin{remark}[Credit is spent only through negative rate]\label{sh:rem:monot}
$D$ is nondecreasing exactly while $\mu\ge0$. Accumulated credit cannot be
destroyed except by first driving the improvement rate negative; this
monotone coupling is the source of all horizon intuition, and its
quantitative teeth are extracted in \S\ref{sh:subsub:softhorizon}.
\end{remark}

\begin{remark}[Nonparametric-Bayes reading of the model]\label{sh:rem:iwp}
Definition~\ref{sh:def:model} places an integrated Wiener process prior on
$-\log h$ --- equivalently, the classical cubic smoothing-spline prior of
Wahba \cite{Wah78}, since a prior whose second derivative is white noise
is precisely one whose first derivative is Brownian. The model is
therefore the canonical smoothness prior on the log-hazard rather than an
exotic construction, and the results below may be read as the survival
asymptotics that prior implies.
\end{remark}

\begin{lemma}[Exact gauge role of the initial hazard]\label{sh:lem:hgauge}
Write $\Lambda^{(h_0)}$ for the cumulative hazard with initial value $h_0$.
Then $\Lambda^{(h_0)}_t=h_0\Lambda^{(1)}_t$ pathwise; hence $\Imm$ does not
depend on $h_0$, while
$\PP(T=\infty)=\E[e^{-h_0\Lambda^{(1)}_\infty}\one_\Imm]$ is continuous and
nonincreasing in $h_0$, strictly decreasing whenever $\PP(\Imm)>0$, with
$\lim_{h_0\downarrow0}\PP(T=\infty)=\PP(\Imm)$ and
$\lim_{h_0\uparrow\infty}\PP(T=\infty)=0$.
\end{lemma}

\begin{proof}
The identity is \eqref{sh:eq:model}; the limits follow by monotone and
dominated convergence.
\end{proof}

Thus $h_0$ can never affect a dichotomy, only a value: it is the first and
most complete gauge degree of freedom. The same will be shown for
$(\mu_0,D_0)$ at the level of asymptotic exponents.

\subsubsection{Dimensional structure, self-similarity, and the gauge orbit}
\label{sh:subsub:dimensional}

\begin{definition}[Units and intrinsic scales]\label{sh:def:units}
Assigning $[t]=\mathsf T$ gives $[\mu]=\mathsf T^{-1}$,
$[\sigma]=\mathsf T^{-3/2}$, $[D]=1$, $[h]=\mathsf T^{-1}$. The unique
timescale constructible from $\sigma$ alone is $\sigma^{-2/3}$; the
time-equivalents of the initial data are
\begin{equation}\label{sh:eq:timescales}
t_\mu:=\frac{\mu_0^2}{\sigma^2},\qquad
t_D:=\Bigl(\frac{D_0}{\sigma}\Bigr)^{2/3}\quad(D_0>0);
\end{equation}
and the unique dimensionless state coordinate on the quadrant
$\{\mu>0,D>0\}$ is the \emph{horizon coordinate}
\begin{equation}\label{sh:eq:kappa}
\kappa:=\frac{\sigma^2D}{\mu^3},\qquad\text{satisfying}\qquad
\frac{t_D}{t_\mu}=\kappa_0^{2/3}.
\end{equation}
Hatted variables are $\hat\mu=\mu\sigma^{-2/3}$, $\hat D=D$,
$\hat h=h\sigma^{-2/3}$, $\hat t=t\sigma^{2/3}$.
\end{definition}

\begin{lemma}[Second-order structure]\label{sh:lem:secondorder}
$\Var(I_t)=t^3/3$, $\Cov(W_t,I_t)=t^2/2$, and for $u\ge1$,
$\E[I_1I_u]=\tfrac u2-\tfrac16$.
\end{lemma}

\begin{proof}
By Fubini, $\E[I_aI_b]=\int_0^a\!\int_0^b(\alpha\wedge\beta)\,d\beta\,d\alpha$.
For $a=b=t$ this is $2\int_0^t\!\int_0^\beta\alpha\,d\alpha\,d\beta=t^3/3$.
Next, $\E[W_tI_t]=\int_0^t(\alpha\wedge t)\,d\alpha=t^2/2$. Finally, for
$u\ge1$, $\int_0^u(\alpha\wedge\beta)\,d\beta
=\alpha u-\alpha^2/2$ for $\alpha\le1\le u$, and integrating over
$\alpha\in[0,1]$ gives $u/2-1/6$.
\end{proof}

\begin{proposition}[Exact self-similarity]\label{sh:prop:selfsim}
Fix $\sigma$ and $c>0$. If $(\mu,D)$ starts from $(\mu_0,D_0)$, then
\[
\bigl(\mu_{c^2t},\,D_{c^2t}\bigr)_{t\ge0}\ \eqd\
\bigl(c\,\mu'_t,\,c^3D'_t\bigr)_{t\ge0},
\]
where $(\mu',D')$ starts from $(\mu_0/c,\,D_0/c^3)$. Equivalently,
$(I_{ct})_{t\ge0}\eqd(c^{3/2}I_t)_{t\ge0}$.
\end{proposition}

\begin{proof}
Brownian scaling $(W_{c^2t})_t\eqd(cW_t)_t$, integrated once in time.
\end{proof}

\begin{proposition}[The gauge orbit]\label{sh:prop:gauge}
The initial data $(h_0,\mu_0,D_0)$ affect asymptotics only through
prefactors and timescales, never through exponents. Precisely:
\emph{(i)} $h_0$ is exactly gauge in the sense of
Lemma~\ref{sh:lem:hgauge};
\emph{(ii)} the deterministic contribution
$x_0-\mu_0t$ to $x_t$ is affine in $t$, hence
$(x_0-\mu_0t)/t^{3/2}\to0$, while $\sigma I_t/t^{3/2}\eqd\sigma I_1$ is
nondegenerate for every $t$: the initial data live in the subspace that the
$t^{3/2}$ fluctuations cannot resolve;
\emph{(iii)} under the scaling of Proposition~\ref{sh:prop:selfsim} the
initial data transform as $(\mu_0,D_0)\mapsto(c\mu_0,c^3D_0)$, so any
scale-invariant asymptotic quantity is constant on the two-parameter orbit
$\{(c\mu_0,c^3D_0):c>0\}$ and, when finite and positive, is a function of
$\kappa_0$ alone.
\end{proposition}

\begin{proof}
(i) is Lemma~\ref{sh:lem:hgauge}; (ii) is immediate from
Proposition~\ref{sh:prop:selfsim} and Lemma~\ref{sh:lem:secondorder};
(iii) holds because the orbits of the scaling action on the open quadrant
are exactly the level curves of $\kappa$.
\end{proof}

The sharp form of the slogan --- that even the \emph{prefactors} of survival
asymptotics depend on the data only through a single equivalent timescale ---
is Proposition~\ref{sh:prop:degrees} below.

\subsubsection{Reduction to a persistence problem}\label{sh:subsub:reduction}

\begin{proposition}[Feynman--Kac]\label{sh:prop:FK}
Let $u(t,x,\mu):=\PP_{x,\mu}(T>t)=\E_{x,\mu}\bigl[e^{-\Lambda_t}\bigr]$.
Then $u$ is the (classical, by Lemma~\ref{sh:lem:hypo}) solution of
\begin{equation}\label{sh:eq:FK}
\partial_tu=\frac{\sigma^2}{2}\partial_\mu^2u-\mu\,\partial_xu-e^{x}u,
\qquad u(0,\cdot,\cdot)\equiv1 .
\end{equation}
\end{proposition}

\begin{proof}
It\^o's formula applied to $s\mapsto
u(t-s,x_s,\mu_s)\exp(-\int_0^se^{x_r}dr)$ exhibits the right side as a
martingale if and only if \eqref{sh:eq:FK} holds; regularity is supplied by
hypoellipticity of the killed generator.
\end{proof}

The killing rate $e^x$ vanishes rapidly for $x\ll0$ and diverges for
$x\gg0$: at exponent level it is an absorbing wall at $x=0$, of thickness
$O(1)$ in $x$ --- a factor $e$ in the hazard --- which is negligible against
excursions of size $t^{3/2}$. We elevate this to a labeled hypothesis, since
it is the single non-rigorous reduction on which the tail asymptotics of
\S\ref{sh:subsub:tail} rest.

\begin{ansatz}[Soft-wall reduction]\label{sh:ansatz}
Let $\tau:=\inf\{t\ge0:D_t\le0\}$ denote the hard-wall passage time of the
Kolmogorov diffusion. Then for every admissible initial datum there exist
constants $0<c\le C<\infty$ and $t_0<\infty$ with
\[
c\,\PP(\tau>t)\ \le\ \PP(T>t)\ \le\ C\,\PP(\tau>t),\qquad t\ge t_0,
\]
after an $O(1)$ adjustment of the initial credit. Equivalently,
$\log\PP(T>t)=\log\PP(\tau>t)+O(1)$.
\end{ansatz}

\begin{remark}[Status of the Ansatz]\label{sh:rem:ansatzstatus}
One half is rigorous modulo a cited input. \emph{Lower bound:} on the event
$\{D_u\ge f(u)\ \forall u\le t\}$ with $f(u)=2\log\frac{2+u}{2}$ one has
$\Lambda_t\le h_0e^{-D_0}\int_0^\infty\bigl(\tfrac{2+u}{2}\bigr)^{-2}du<\infty$,
whence $\PP(T>t)\ge e^{-C}\,\PP(D\ge f\text{ on }[0,t])$; that a logarithmic
moving boundary does not change the persistence exponent of an integrated
process is the standard moving-boundary phenomenon (see the survey
\cite{AS15} and references therein; cf.\ also the robustness expressed
by the universality results \cite{AD13,DW15}). \emph{Upper bound:} from
$\PP(T>t)\le e^{-K}+\PP(\Lambda_t\le K)$ and
$\Lambda_t\ge h_0e^{-D_0}\Leb\{u\le t:D_u\le0\}$, the event
$\{\Lambda_t\le K\}$ forces $D$ to remain positive off a set of bounded
Lebesgue measure. What is missing is the (plausible, unproved here)
statement that \emph{persistence with a bounded occupation tolerance has the
same exponent as strict persistence}. Everything tagged \textbf{[A]} below
depends on exactly this gap and nothing else; \S\ref{sh:sub:numerics}
describes its direct numerical test.
\end{remark}

% ---------------------------------------------------------------------
\subsection{The driftless verdict: certain death, infinite expectation}
\label{sh:sub:verdict}

\subsubsection{Impossibility of immortality}\label{sh:subsub:noimm}

\begin{theorem}[No viability; \textbf{P}]\label{sh:thm:noimm}
For every $(h_0,\mu_0,D_0,\sigma)$ in the driftless model
\eqref{sh:eq:model}, $\Lambda_\infty=\infty$ almost surely. Consequently
$S_\infty=0$ a.s., $\PP(\Imm)=0$, $\PP(T=\infty)=0$, and $T<\infty$ a.s.
\end{theorem}

\begin{proof}
Since $\Lambda_\infty=h_0e^{-D_0}\int_0^\infty e^{-\mu_0u-\sigma I_u}du$
and the prefactor is a fixed positive constant, it suffices to prove that
the integral diverges a.s.

\emph{Step 1 (Lamperti transform).} Define $Z_s:=e^{-3s/2}I_{e^s}$,
$s\ge0$. As a linear functional of $W$, $Z$ is a centered Gaussian process.
By Proposition~\ref{sh:prop:selfsim}, for each $r\ge0$,
$(Z_{s+r})_{s\ge0}=(e^{-3(s+r)/2}I_{e^re^s})_{s\ge0}
\eqd(e^{-3s/2}I_{e^s})_{s\ge0}=(Z_s)_{s\ge0}$: $Z$ is stationary. By
Lemma~\ref{sh:lem:secondorder}, for $s\ge0$,
\begin{equation}\label{sh:eq:lampcov}
\E[Z_0Z_s]=e^{-3s/2}\Bigl(\frac{e^{s}}{2}-\frac16\Bigr)
=\frac12e^{-s/2}-\frac16e^{-3s/2}\ \longrightarrow\ 0 .
\end{equation}

\emph{Step 2 (ergodicity).} A stationary Gaussian process whose covariance
vanishes at infinity is mixing \cite{Mar49}, hence ergodic.

\emph{Step 3 (Birkhoff).} Fix $\varepsilon>0$. Since
$Z_0\sim N(0,\tfrac13)$, $p_\varepsilon:=\PP(Z_0\le-\varepsilon)>0$. By the
continuous-time Birkhoff theorem,
$\tfrac1S\Leb\{s\in[0,S]:Z_s\le-\varepsilon\}\to p_\varepsilon$ a.s., so
$\Leb\{s\ge0:Z_s\le-\varepsilon\}=\infty$ a.s.

\emph{Step 4 (change of variables).} Let
$A:=\{u\ge1:I_u\le-\varepsilon u^{3/2}\}$. Under $u=e^s$ the set in Step~3
maps into $A$ with
$\Leb\{s:Z_s\le-\varepsilon\}=\int_A\frac{du}{u}\le\Leb(A)$,
the inequality because $u\ge1$. Hence $\Leb(A)=\infty$ a.s.

\emph{Step 5 (divergence).} Choose $u_1=u_1(\varepsilon,\mu_0,\sigma)$ so
that $\sigma\varepsilon u^{3/2}-\mu_0u\ge0$ for $u\ge u_1$. Then
\[
\int_0^\infty e^{-\mu_0u-\sigma I_u}\,du
\ \ge\ \int_{A\cap[u_1,\infty)}e^{\sigma\varepsilon u^{3/2}-\mu_0u}\,du
\ \ge\ \Leb\bigl(A\cap[u_1,\infty)\bigr)\ =\ \infty
\qquad\text{a.s.}\qedhere
\]
\end{proof}

\begin{remark}[The asymmetry of exponential response]\label{sh:rem:asym}
The driver is symmetric; the response is not. A favorable excursion
suppresses the integrand $e^{-D}$ multiplicatively toward zero ---
\emph{bounded} benefit for $\Lambda$ --- while an unfavorable excursion of
relative size $\varepsilon$ inflates it without bound. Exponential response
to a symmetric driver is structurally hostile to survival; the mechanism is
the sign-oscillation that makes low-dimensional random walks recurrent,
aggravated by this asymmetry. Identifying exactly what a horizon must break
--- the symmetry of the driver, not the response --- is the burden of
\S\ref{sh:subsub:scale}.
\end{remark}

\subsubsection{The survival tail and its consequences}\label{sh:subsub:tail}

\begin{theorem}[Persistence of integrated Brownian motion; \textbf{C}]
\label{sh:thm:persist}
Let $(\mu,D)$ start from $(m,0)$ with $m>0$ and let
$\tau=\inf\{t:D_t\le0\}$. There is a constant $c_1>0$ such that
\begin{equation}\label{sh:eq:persist}
\PP_{(m,0)}(\tau>t)\ \sim\ c_1\Bigl(\frac{m^2}{\sigma^2t}\Bigr)^{1/4},
\qquad t\to\infty .
\end{equation}
The exponent $\tfrac14$ --- half the Brownian value, halved by one
integration --- is the persistence exponent of integrated Brownian
motion, and its provenance is a chain worth recording precisely. McKean
\cite{McK63} computed the first-passage (``winding'') law of the
Kolmogorov diffusion; Goldman \cite{Gol71} derived the passage-rate
asymptotic from McKean's law; Lachal \cite{Lac91} obtained the exact
first-passage distributions; Sinai \cite{Sin92} initiated and settled,
up to constants, the discrete counterpart $n^{-1/4}$ for the integrated
random walk; Isozaki--Watanabe \cite{IW94} supplied the sharp constants;
and Denisov--Wachtel \cite{DW15} proved the exact $n^{-1/4}$ asymptotics
universally for centered integrated random walks with $2+\delta$
moments, via a discrete potential theory (see also the universality
result of Aurzada--Dereich \cite{AD13} for integrated processes).
Starting states in the open quadrant obey the same power with a
prefactor proportional to the harmonic function of
Proposition~\ref{sh:prop:harmonic}; this general-start form is
established in \cite{GJW99}, which recovers \cite{Gol71,IW94} as
special cases. For the surrounding theory see the survey \cite{AS15}
and, on the physics side --- where the model is the \emph{random
acceleration process} and has an extensive literature of its own ---
the review \cite{BMS13}.
\end{theorem}

\begin{remark}[The $m^{1/2}$ law is exact given the $m=1$ asymptotic]
By Proposition~\ref{sh:prop:selfsim} with $c=1/m$ (and $\sigma=1$),
$\PP_{(m,0)}(\tau>t)=\PP_{(1,0)}(\tau>t/m^2)$, so once
\eqref{sh:eq:persist} is known for one $m$ it holds for all, prefactor
scaling as $m^{1/2}$ exactly. The initial rate enters through its
time-equivalent $t_\mu$ and nothing else.
\end{remark}

\begin{claim}[Survival tail; \textbf{A}]\label{sh:claim:tail}
Under Ansatz~\ref{sh:ansatz}, for every initial datum there is
$C=C(\hat h_0,\hat\mu_0,\hat D_0)\in(0,\infty)$ with
\begin{equation}\label{sh:eq:tail}
\PP(T>t)\ \asymp\ C\,t^{-1/4},\qquad t\to\infty .
\end{equation}
\end{claim}

All of the following consequences are elementary integrations of
\eqref{sh:eq:tail}; each inherits the tag \textbf{[A]}. Write
$\alpha=\tfrac14$.

\begin{corollary}[Moments]\label{sh:cor:moments}
$\E[T^p]<\infty$ if and only if $p<\tfrac14$. In particular every moment of
order at least $\tfrac14$ diverges: $\E[\sqrt T]=\E[T^{1/3}]=\infty$; the
law of $T$ is heavier-tailed than Cauchy.
\end{corollary}

\begin{proof}
$\E[T^p]=\int_0^\infty pt^{p-1}\PP(T>t)\,dt$ converges at infinity iff
$p-1-\tfrac14<-1$.
\end{proof}

\begin{corollary}[Infinite life expectancy, for every initial condition]
\label{sh:cor:mean}
$\E[T]=\infty$ while $\PP(T<\infty)=1$, for all
$(h_0,\mu_0,D_0,\sigma)$. This is the null-recurrent signature: the
two-dimensional random-walk regime (certain return, infinite expected
return time), not the three-dimensional one. Initial conditions cannot
rescue the mean because, by Proposition~\ref{sh:prop:gauge}, they enter
\eqref{sh:eq:tail} only through the prefactor.
\end{corollary}

\begin{corollary}[Truncated mean; observational artifact]\label{sh:cor:trunc}
For an observation deadline $\mathcal T$,
$\E[T\wedge\mathcal T]=\int_0^{\mathcal T}\PP(T>t)\,dt
\asymp\tfrac{4C}{3}\,\mathcal T^{3/4}$: any finite dataset reports a finite
mean lifetime that is an artifact of its own duration, growing as the
$3/4$ power of the window and never converging.
\end{corollary}

\begin{corollary}[Quantiles]\label{sh:cor:quant}
If $t_q$ solves $\PP(T>t_q)=q$ then $t_q\asymp(C/q)^4$: lifetime quantiles
diverge as the \emph{fourth power} of rarity; the $1$-percent-longest
lifetime exceeds the $10$-percent-longest by a factor of order $10^4$.
\end{corollary}

\begin{corollary}[Residual life is Pareto]\label{sh:cor:residual}
For $\lambda\ge1$,
$\PP(T>\lambda t\mid T>t)\to\lambda^{-1/4}$: residual lifetime, in units of
attained age, converges to a Pareto law of index $\tfrac14$ ---
scale-free, with infinite mean at every age. Under a deadline
$\mathcal T$,
\[
\E\bigl[(T\wedge\mathcal T)-t\,\big|\,T>t\bigr]
=\int_t^{\mathcal T}\Bigl(\frac ut\Bigr)^{-1/4}du
=\frac43\,t^{1/4}\bigl(\mathcal T^{3/4}-t^{3/4}\bigr).
\]
\end{corollary}

\begin{corollary}[Effective macroscopic hazard; Lindy law]\label{sh:cor:lindy}
The effective cumulative hazard of the death time is
$-\log\PP(T>t)=\tfrac14\log t+O(1)$; at the level of the asymptotics the
effective hazard rate is $r(t)\approx\tfrac1{4t}$. Since
$\int^\infty\tfrac{dt}{4t}=\infty$, death remains certain; immortality
would require an integrable effective hazard. The system thus sits at tail
exponent $\tfrac14$: far from the mortality boundary (any positive exponent
kills) and deep inside the infinite-mean regime (boundary at exponent $1$).
\end{corollary}

\subsubsection{The conditioning ladder}\label{sh:subsub:ladder}

We quantify how each conditioning enhances survival.

\begin{proposition}[Degrees of the initial data; \textbf{C/A}]
\label{sh:prop:degrees}
\emph{(i) Initial rate.} From $(m,0)$ with $m>0$: the rate stays positive
until $W$ reaches $-m/\sigma$, a time of order $t_\mu=m^2/\sigma^2$ during
which $D$ increases; thereafter, exactly by scaling,
$\PP_{(m,0)}(\tau>t)\sim c_1(t_\mu/t)^{1/4}\propto m^{1/2}$
\emph{(degree $\tfrac12$: a tenfold gain in survival odds costs a
hundredfold increase in $m$)}.
\emph{(ii) Initial credit.} From $(0,d)$ with $d>0$: by scaling with
$c=d^{-1/3}$, $\PP_{(0,d)}(\tau>t)=\PP_{(0,1)}(\sigma^{2/3}\tau>
t\,d^{-2/3}\sigma^{2/3})\asymp(t_D/t)^{1/4}\propto d^{1/6}$
\emph{(degree $\tfrac16$)}.
\emph{(iii) Unification.} All initial data act only through the equivalent
timescale they generate:
\begin{equation}\label{sh:eq:tstar}
\PP(T>t)\ \asymp\ \Bigl(\frac{t_*}{t}\Bigr)^{1/4},\qquad
t_*\ \asymp\ \sigma^{-2/3}\max\bigl(1,\ \hat\mu_0^{\,2},\ \hat
D_0^{\,2/3}\bigr),
\end{equation}
and which term dominates is decided by the horizon coordinate: by
\eqref{sh:eq:kappa}, $t_D/t_\mu=\kappa_0^{2/3}$, so rate dominates for
$\kappa_0\ll1$ and credit for $\kappa_0\gg1$.
\end{proposition}

\begin{proof}
The scaling identities are exact (Proposition~\ref{sh:prop:selfsim});
transfer to $\PP(T>t)$ is Ansatz~\ref{sh:ansatz}. That the exponent from
interior starting states is unchanged is part of
Theorem~\ref{sh:thm:persist}.
\end{proof}

\begin{proposition}[Homogeneity of the harmonic function; \textbf{P/C}]
\label{sh:prop:harmonic}
On the open quadrant define
\[
\hm(\mu,d)\ :=\ \lim_{t\to\infty}t^{1/4}\,\PP_{(\mu,d)}(\tau>t);
\]
the limit exists, is finite and positive \cite{IW94,GJW99}. Then $\hm$ is
invariant-harmonic for the killed Kolmogorov diffusion and homogeneous of
degree $\tfrac12$ under the scaling of
Proposition~\ref{sh:prop:selfsim}, $\hm(c\mu,c^3d)=c^{1/2}\hm(\mu,d)$.
Consequently
\begin{equation}\label{sh:eq:harm}
\hm(\mu,d)=\mu^{1/2}F\bigl(\kappa\bigr)=d^{1/6}\tilde F(\kappa),
\qquad \kappa=\sigma^2d/\mu^3,
\end{equation}
for profiles $F,\tilde F$ determined by the explicit
confluent-hypergeometric formulas constructed in \cite{GJW99}. (We flag
a provenance point: the special functions arising here are confluent
hypergeometric, not Airy; Airy functions belong to the neighbouring
parabolic-drift literature.)
\end{proposition}

\begin{proof}
Given existence of the limit, homogeneity is immediate from
Proposition~\ref{sh:prop:selfsim}:
\[
t^{1/4}\,\PP_{(c\mu,c^3d)}(\tau>t)
= c^{1/2}\,\bigl(t/c^2\bigr)^{1/4}\,\PP_{(\mu,d)}\bigl(\tau>t/c^{2}\bigr)
\ \longrightarrow\ c^{1/2}\,\hm(\mu,d).
\]
Harmonicity follows from the Markov property and the same limit taken one
step later. Representation \eqref{sh:eq:harm} is homogeneity restricted to
the orbit coordinates of Proposition~\ref{sh:prop:gauge}(iii).
\end{proof}

\begin{remark}[Conditioning on survival forever: the Doob transform;
\textbf{C}]\label{sh:rem:doob}
Conditioning on $\{T=\infty\}$ is vacuous (Theorem~\ref{sh:thm:noimm});
its canonical regularization is the $\hm$-transform, the
$t\to\infty$ limit of conditioning on $\{\tau>t\}$. Under the transformed
measure the process never returns to zero credit and, by self-similarity,
\[
\mu_s\ \asymp\ \sigma s^{1/2},\qquad D_s\ \asymp\ \sigma s^{3/2},
\qquad\text{hence}\qquad \log h_s\ \sim\ -\sigma s^{3/2}:
\]
conditioned survivors exhibit superexponentially decaying hazard and
appear, from the inside, unambiguously immortal --- although the
unconditioned process is certainly mortal and the model contains no drift.
The drift is manufactured by selection; this is the strongest true form of
the event-horizon intuition, and it is a statement about ensembles, not
about state space. The accompanying Yaglom-type statement --- convergence
in law of $(\mu_t/\sigma t^{1/2},\,D_t/\sigma t^{3/2})$ conditionally on
$\{\tau>t\}$ to a nondegenerate limit --- is the correct
``stationary profile of a long-lived system,'' stationary only in
self-similar coordinates; see \cite{GJW99,AS15}, and \cite{MV12} for
the general theory of quasi-stationary limits against which this
self-similar variant should be read.
Section~\ref{sh:sub:interp} promotes this remark to the section's central
result: the conditioning that manufactures the drift becomes itself the
object of study, and the transform described here is identified as the
exact limit of the survivorship interpolation.
\end{remark}

% ---------------------------------------------------------------------
\subsection{Horizon theory}\label{sh:sub:horizon}

\subsubsection{The driftless horizon is a timescale, and it is maximally
soft}\label{sh:subsub:softhorizon}

In the driftless model the naive horizon quantity
$\Pi(\mu,D):=\PP_{(\mu,D)}(D\text{ never returns to }0)$ vanishes
identically (Theorem~\ref{sh:thm:noimm}): null recurrence leaves no escape
probability to be gradual about, and the horizon must be recast in the only
currency the model possesses --- time.

\begin{proposition}[Recovery cost is sublinear in credit; \textbf{P}]
\label{sh:prop:recovery}
Let $\rho(d)$ denote the median of the return time
$\inf\{t:D_t\le0\}$ from $(\mu,D)=(0,d)$, $d>0$. Then exactly
\begin{equation}\label{sh:eq:recovery}
\rho(d)=\rho(1)\,\Bigl(\frac d\sigma\Bigr)^{2/3}\sigma^{2/3}\cdot
\sigma^{-2/3}=\rho_1\Bigl(\frac d\sigma\Bigr)^{2/3},
\qquad \rho_1:=\rho(1)\big|_{\sigma=1}\in(0,\infty),
\end{equation}
and the return time is finite a.s.\ with tail index $\tfrac14$. Doubling
accumulated credit multiplies the recovery timescale by only
$2^{2/3}\approx1.59$: credit is far easier to undo than a linear intuition
suggests.
\end{proposition}

\begin{proof}
Finiteness a.s.\ is contained in the proof of
Theorem~\ref{sh:thm:noimm} (the set $A$ there is unbounded). The scaling
identity is Proposition~\ref{sh:prop:selfsim} with $c=d^{-1/3}$ applied to
the median, a scale-equivariant functional; positivity and finiteness of
$\rho_1$ follow since the return time is a nondegenerate positive random
variable. The tail index is Theorem~\ref{sh:thm:persist} read for the
returning coordinate.
\end{proof}

\begin{proposition}[Horizon level sets are $\kappa$-rays; \textbf{P}]
\label{sh:prop:rays}
Every scale-invariant function on the open quadrant is constant exactly on
the curves $D\propto\mu^3/\sigma^2$; if a horizon exists in any
scale-invariant sense, it is a ray of constant $\kappa$, and $\kappa$ is
its natural transverse coordinate.
\end{proposition}

\begin{proof}
Proposition~\ref{sh:prop:gauge}(iii).
\end{proof}

\begin{definition}[Deadline-relative horizon and its width]
\label{sh:def:deadline}
Fix a time budget $\mathcal T$. The \emph{deadline horizon} is the level
set $\{t_*(\mu,D)\asymp\mathcal T\}$, located at
$\mu^*\asymp\sigma\mathcal T^{1/2}$, $D^*\asymp\sigma\mathcal T^{3/2}$.
Its \emph{softness} is the logarithmic sensitivity of survival to state.
\end{definition}

\begin{corollary}[Maximal softness; \textbf{A}]\label{sh:cor:soft}
Under Claim~\ref{sh:claim:tail},
$\dfrac{\partial\log\PP(T>\mathcal T)}{\partial\log t_*}=\dfrac14+o(1)$.
Changing survival probability by a factor $e$ therefore requires changing
$t_*$ by $e^4\approx55$ and, since $t_*\propto D^{2/3}$ in the
credit-dominated regime, changing $D$ by $e^6\approx403$. The driftless
horizon is smeared over roughly two orders of magnitude in timescale and
nearly three in credit: there is no regime in which resistance has
arbitrarily little effect, only a $\tfrac14$-power degradation --- about as
gentle a degradation as a power law permits.
\end{corollary}

\subsubsection{The scale-function criterion: P\'olya's theorem, correctly
translated}\label{sh:subsub:scale}

Let the driver be a regular diffusion on $\R$,
\begin{equation}\label{sh:eq:generaldriver}
d\mu_t=b(\mu_t)\,dt+\varsigma(\mu_t)\,dW_t,\qquad
s(\mu):=\int_{0}^{\mu}\exp\Bigl(-\int_{0}^{v}\frac{2b(w)}{\varsigma^2(w)}
\,dw\Bigr)dv,
\end{equation}
with $\varsigma>0$ locally bounded away from $0$ and $\infty$; $s$ is the
scale function and Feller's test (see \cite[Ch.~15]{KT81} or
\cite[Ch.~VII]{RY99}) classifies the boundary behavior.

\begin{theorem}[Trichotomy; \textbf{P}]\label{sh:thm:trichotomy}
For the model \eqref{sh:eq:model} with driver
\eqref{sh:eq:generaldriver}:
\begin{enumerate}[label=\emph{(\roman*)},itemsep=1pt]
\item If $s(+\infty)<\infty$ and $s(-\infty)=-\infty$, then
$\mu_t\to+\infty$ a.s., $\Lambda_\infty<\infty$ a.s., $\PP(\Imm)=1$, and
$\PP(T=\infty)=\E[e^{-\Lambda_\infty}]\in(0,1)$.
\item If $s(-\infty)>-\infty$ and $s(+\infty)=\infty$, then
$\mu_t\to-\infty$ a.s., $\Lambda_\infty=\infty$ a.s., and $T<\infty$ a.s.
\item If both limits are finite, then
\begin{equation}\label{sh:eq:scaleformula}
\PP(\Imm)\;=\;\frac{s(\mu_0)-s(-\infty)}{s(+\infty)-s(-\infty)},
\end{equation}
with $\Lambda_\infty<\infty$ on $\{\mu_t\to+\infty\}$ and
$\Lambda_\infty=\infty$ on the complementary event
$\{\mu_t\to-\infty\}$.
\item If both limits are infinite, the driver is recurrent and the outcome
is \emph{not} determined by $s$; see
Proposition~\ref{sh:prop:recurrent}.
\end{enumerate}
\end{theorem}

\begin{proof}
(i) $\PP(\lim_t\mu_t=+\infty)=1$ is Feller's test \cite{KT81}. On this
event choose $t_0(\omega)$ with $\mu_t\ge1$ for $t\ge t_0$; then
$D_t\ge D_{t_0}+(t-t_0)$, so
$\Lambda_\infty\le\Lambda_{t_0}+h_0e^{-(D_{t_0}-D_0)}\int_0^\infty
e^{-v}dv<\infty$ a.s. Positivity of $\PP(T=\infty)$ follows since
$e^{-\Lambda_\infty}>0$ a.s.; it is $<1$ since $\Lambda_\infty\ge
\Lambda_1>0$ a.s. (ii) is the mirror image: eventually $\mu_t\le-1$, so
$e^{-(D_t-D_0)}\to\infty$ and $\Lambda_\infty=\infty$; death then occurs in
finite time a.s.\ by Definition~\ref{sh:def:cox}. (iii) The hitting
probabilities of $\pm\infty$ are the classical scale-function formula, and
the two events partition the space up to null sets; apply the pathwise
arguments of (i) and (ii) on each. (iv) Recurrence under unbounded scale is
Feller's test again; Example~\ref{sh:ex:recurrentcounter} shows $s$ alone
cannot decide.
\end{proof}

\begin{proposition}[Recurrent drivers; \textbf{P}]\label{sh:prop:recurrent}
Suppose the driver \eqref{sh:eq:generaldriver} is positive recurrent with
stationary law $\pi$ and $\int|\mu|\,\pi(d\mu)<\infty$; write
$\bar m:=\int\mu\,\pi(d\mu)$.
\begin{enumerate}[label=\emph{(\alph*)},itemsep=1pt]
\item If $\bar m>0$ then $D_t/t\to\bar m$ a.s., $\Lambda_\infty<\infty$
a.s., and $\PP(\Imm)=1$.
\item If $\bar m<0$ then $\Lambda_\infty=\infty$ a.s.\ and $T<\infty$ a.s.
\item If $\bar m=0$ and the cycle integral defined in the proof is
nondegenerate (automatic for nondegenerate diffusions), then
$\Lambda_\infty=\infty$ a.s.: \emph{the critical case falls on the mortal
side.}
\end{enumerate}
For null-recurrent drivers no general statement is offered
\textbf{[O]}; the Brownian driver, the fundamental null-recurrent case, is
settled unconditionally by Theorem~\ref{sh:thm:noimm}.
\end{proposition}

\begin{proof}
(a), (b): the ergodic theorem for positive recurrent one-dimensional
diffusions gives $D_t/t\to\bar m$ a.s.\ (see \cite[Ch.~X]{RY99}); if
$\bar m>0$ then eventually $D_t\ge D_0+\tfrac{\bar m}2t$ and
$\Lambda_\infty<\infty$; if $\bar m<0$ the integrand
$e^{-(D_t-D_0)}\to\infty$.

(c) Fix interior levels $a<b$ and regenerate at up-crossings: let
$R_0:=\inf\{t:\mu_t=b\}$ and inductively
$S_k:=\inf\{t>R_{k-1}:\mu_t=a\}$, $R_k:=\inf\{t>S_k:\mu_t=b\}$. By the
strong Markov property the cycles over $[R_{k-1},R_k)$ are i.i.d.; positive
recurrence gives $\E[L]<\infty$ for the cycle length
$L:=R_1-R_0$, and the cycle occupation identity yields, for the cycle
integral $\xi_k:=\int_{R_{k-1}}^{R_k}\mu_s\,ds$,
\[
\E[\xi_1]=\E[L]\int\mu\,d\pi=0,\qquad
\E|\xi_1|\le\E\Bigl[\int_{\text{cycle}}|\mu_s|ds\Bigr]
=\E[L]\int|\mu|\,d\pi<\infty .
\]
Thus $S_n:=D_{R_n}-D_{R_0}=\sum_{k\le n}\xi_k$ is a nondegenerate mean-zero
i.i.d.\ walk, so $\liminf_nS_n=-\infty$ a.s.\ by the classical
oscillation dichotomy for random walks (see, e.g., \cite{Dur19});
in particular the events $B_n:=\{D_{R_n}\le-1\}$ occur infinitely often
a.s. Choose $\delta>0$ and $M<\infty$ with
\[
q:=\PP_b\Bigl(L\ge\delta,\ \sup_{0\le u\le\delta}
\Bigl|\int_0^u\mu_s\,ds\Bigr|\le M\Bigr)>0,
\]
possible since $L>0$ a.s.\ and the supremum is finite a.s. Let
$G_n$ be the corresponding event for cycle $n{+}1$; by the strong Markov
property $\PP(G_n\mid\Fc_{R_n})=q$, so
$\sum_n\PP(B_n\cap G_n\mid\Fc_{R_n})=q\sum_n\one_{B_n}=\infty$ a.s., and
L\'evy's extension of the Borel--Cantelli lemma (see, e.g.,
\cite{Dur19}) gives
$B_n\cap G_n$ infinitely often a.s. On $B_n\cap G_n$, for
$u\in[0,\delta]$ one has $D_{R_n+u}\le-1+M$, hence
\[
\int_{R_n}^{R_n+\delta}e^{-(D_s-D_0)}\,ds\ \ge\ \delta\,e^{\,1-M+D_0}\,>\,0 ,
\]
and these windows are disjoint. Summing over infinitely many of them,
$\Lambda_\infty=\infty$ a.s.
\end{proof}

\begin{example}[Recurrence alone does not decide]\label{sh:ex:recurrentcounter}
The mean-reverting Ornstein--Uhlenbeck driver
$d\mu=\theta(\bar\mu-\mu)dt+\sigma dW$ with $\theta>0$ has Gaussian-growing
$s'$, hence $s(\pm\infty)=\pm\infty$: it is (positive) recurrent for every
$\bar\mu$. Yet Proposition~\ref{sh:prop:recurrent} gives
$\PP(\Imm)=1$ for $\bar\mu>0$ and certain death for $\bar\mu\le0$. A
horizon exists here --- sharp, at $\bar\mu=0$ --- but it lives in
\emph{parameter} space, not in the initial condition, which is again pure
gauge.
\end{example}

\begin{remark}[Two warnings]\label{sh:rem:warnings}
Two traps in this circle of ideas deserve explicit flags.
\emph{First}, it is tempting to summarize \eqref{sh:eq:scaleformula} as
``recurrent driver $\Rightarrow$ certain death''; Example
\ref{sh:ex:recurrentcounter} shows this is false --- in the recurrent class
the order parameter is the stationary mean $\bar m$, with the critical case
$\bar m=0$ mortal, and only in the null-recurrent Brownian case is
symmetry itself the whole story.
\emph{Second}, under adverse drift it is tempting to read the survival tail
off the typical path, which yields doubly-exponential decay; in fact the
tail is only \emph{exponential}, dominated by an optimal ``fighting-back''
fluctuation --- Proposition~\ref{sh:prop:advdrift} below --- a
Freidlin--Wentzell phenomenon in miniature.
\end{remark}

\begin{center}
\begin{tabular}{@{}ll@{}}
\toprule
\textbf{Random walk (P\'olya)} & \textbf{Hazard dynamics}\\
\midrule
dimension $d$ & scale function $s$ of the driver\\
$d\le2$: recurrent, no escape & $s(\pm\infty)=\pm\infty$: no viability from
data\\
$d\ge3$: escape probability $>0$ & $s(\pm\infty)$ finite: viability
probability in $(0,1)$\\
$d=2$: null recurrent & Brownian driver: certain death, infinite mean
life\\
escape probability from $x$ & normalized scale function at $\mu_0$\\
\bottomrule
\end{tabular}
\end{center}

\begin{remark}[Slogan]\label{sh:rem:slogan}
\emph{Existence} of a horizon in the initial condition
$=$ boundedness of the scale function;
\emph{sharpness} of the horizon $=$ steepness of the scale function.
The Brownian driver has $s(\mu)=\mu$, unbounded and of constant slope: its
horizon degenerates not by becoming infinitely sharp but by becoming
infinitely smeared --- a linear scale function distinguishes no location.
Finally, recall from \eqref{sh:eq:twotier} that
\eqref{sh:eq:scaleformula} computes \emph{viability}; immortality proper is
$\E[e^{-\Lambda_\infty}\one_\Imm]$, strictly smaller, and this is the
unique juncture at which the gauge parameter $h_0$ re-enters
(Lemma~\ref{sh:lem:hgauge}) --- affecting values, never dichotomies.
\end{remark}

\subsubsection{An exactly solvable gradual horizon}\label{sh:subsub:solvable}

\begin{theorem}[Unstable Ornstein--Uhlenbeck driver; \textbf{P}]
\label{sh:thm:uou}
Let the driver be
\begin{equation}\label{sh:eq:uou}
d\mu_t=\nu\,\mu_t\,dt+\sigma\,dW_t,\qquad \nu>0 .
\end{equation}
Then $e^{-\nu t}\mu_t\to Z$ a.s.\ and in $L^2$, where
$Z\sim N\!\bigl(\mu_0,\ \sigma^2/2\nu\bigr)$, and
\begin{equation}\label{sh:eq:sigmoid}
\PP(\Imm)\;=\;\PP(Z>0)\;=\;\Phi\!\Bigl(\frac{\mu_0\sqrt{2\nu}}{\sigma}\Bigr).
\end{equation}
On $\{Z>0\}$ the hazard decays doubly exponentially,
$\log h_t\sim-(Z/\nu)e^{\nu t}$, and $\Lambda_\infty<\infty$; on
$\{Z<0\}$ the hazard explodes doubly exponentially and $T<\infty$ a.s.
\end{theorem}

\begin{proof}
Variation of constants gives
$\mu_t=e^{\nu t}\bigl(\mu_0+\sigma\int_0^te^{-\nu s}dW_s\bigr)
=:e^{\nu t}M_t$. The martingale $M$ has
$\E\langle M\rangle_\infty=\sigma^2\!\int_0^\infty e^{-2\nu s}ds
=\sigma^2/2\nu<\infty$, hence converges a.s.\ and in $L^2$ to a Gaussian
limit $Z$ with the stated mean and variance. On $\{Z>0\}$ pick
$t_0(\omega)$ with $\mu_t\ge\tfrac Z2e^{\nu t}$ for $t\ge t_0$; then
$D_t\ge D_{t_0}+\tfrac{Z}{2\nu}(e^{\nu t}-e^{\nu t_0})$ and the integrand
$e^{-(D_t-D_0)}$ is dominated by a doubly-exponentially decaying function,
so $\Lambda_\infty<\infty$. On $\{Z<0\}$ symmetrically
$\mu_t\le\tfrac Z2e^{\nu t}$ eventually, $D_t\to-\infty$ doubly
exponentially, the integrand diverges, $\Lambda_\infty=\infty$, and
$T<\infty$ a.s. Finally $\PP(Z=0)=0$, and
$\PP(Z>0)=\Phi(\mu_0/\sqrt{\sigma^2/2\nu})$.
\end{proof}

\begin{remark}[Consistency with the scale-function formula]
\label{sh:rem:crosscheck}
For \eqref{sh:eq:uou}, $s'(v)=\exp(-\nu v^2/\sigma^2)$, so
$s(\pm\infty)$ are finite and the normalized scale function at $\mu_0$ is
the mass below $\mu_0$ of the centered Gaussian with variance
$\sigma^2/2\nu$ --- precisely \eqref{sh:eq:sigmoid}. The exactly solvable
case is thus the generic case of
Theorem~\ref{sh:thm:trichotomy}(iii) in its purest form.
\end{remark}

\begin{definition}[Horizon profile, location, width]\label{sh:def:width}
For a driver with bounded scale, the \emph{horizon profile} is
$\Pi(\mu_0):=\PP(\Imm\mid\mu_0)$, the \emph{location} its median point
$\Pi=\tfrac12$, and the \emph{width} the natural spread of the transition,
e.g.\ the interquartile range of the measure $d\Pi$. For
\eqref{sh:eq:uou}: location $\mu_0=0$; width
\begin{equation}\label{sh:eq:width}
w=\frac{\sigma}{\sqrt{2\nu}},\qquad\text{sharpness}\quad
\Sigma:=w^{-1}=\frac{\sqrt{2\nu}}{\sigma}.
\end{equation}
As $\Sigma\to\infty$ (small noise or strong instability),
$\Pi\to\one\{\mu_0>0\}$: a hard event horizon. As $\Sigma\to0$,
$\Pi\to\tfrac12$ everywhere: no horizon, a fair coin.
\end{definition}

\begin{remark}[Sharpness is a small-noise limit; general principle]
\label{sh:rem:FW}
The width \eqref{sh:eq:width} is exactly the scale of fluctuation of the
terminal charge $Z$ capable of overturning the deterministic verdict
$\operatorname{sign}(\mu_0)$: the sharp horizon is the Freidlin--Wentzell
limit \cite{FW12} of the soft family. We record the principle in the form
answering the question that motivates this section: \emph{a horizon is
never intrinsically sharp; its width is the fluctuation scale of the
driver, and it appears sharp exactly when that scale is small against the
spread of initial conditions.} Operationally, comparing $w$ to the
population spread of $\mu_0$ yields three regimes: $w\ll$ spread --- sharp
horizon, the population splits into the fated; $w\sim$ spread --- genuinely
gradual horizon, initial conditions matter but do not determine;
$w\gg$ spread --- no horizon, the outcome is noise.
\end{remark}

\subsubsection{Model atlas and the role of driver smoothness}
\label{sh:subsub:atlas}

\begin{proposition}[Adverse linear drift: the price of survival;
\textbf{P/S}]\label{sh:prop:advdrift}
Let $\mu_t=\mu_0+\nu t+\sigma W_t$ with $\nu<0$. Then:
\begin{enumerate}[label=\emph{(\roman*)},itemsep=1pt]
\item \textbf{[P]}
$\displaystyle\liminf_{t\to\infty}\frac1t\log\PP(T>t)\ \ge\
-\frac{\nu^2}{2\sigma^2}$;
\item \textbf{[P]} there exists $c>0$ with $\PP(T>t)\le e^{-ct}$ for all
large $t$;
\item \textbf{[S]} the exact rate is
\begin{equation}\label{sh:eq:rate}
\lim_{t\to\infty}\frac1t\log\PP(T>t)\ =\ -\frac{3\nu^2}{8\sigma^2},
\end{equation}
attained by the optimal survival profile
$\mu_{tv}\approx\tfrac{|\nu|t}{4}\,v\,(2-3v)$,
$D_{tv}\approx\tfrac{|\nu|t^2}{4}\,v^2(1-v)$, $v\in[0,1]$: the conditioned
survivor front-loads credit, lets the rate go negative at $v=\tfrac23$, and
coasts so as to exhaust its credit exactly at the deadline.
\end{enumerate}
\end{proposition}

\begin{proof}
(i) By the Markov property and a finite-time event of positive probability
we may assume $\mu_0\ge1$; constants below depend on the data. Let
$a:=|\nu|/\sigma$, $b_0:=(1-\mu_0)/\sigma\le0$, and
\[
A_t:=\Bigl\{W_u\ge b_0+au\ \ \forall u\in[0,t],\ \ W_t\le b_0+at+1\Bigr\}.
\]
On $A_t$, $\mu_u\ge1$ for all $u\le t$, hence $D_u\ge D_0+u$ and
$\Lambda_t\le h_0$, so $\PP(T>t)\ge e^{-h_0}\PP(A_t)$. Girsanov with drift
$a$ (density $\exp(-a\widetilde W_t-a^2t/2)$, $\widetilde W:=W-a\,\cdot$
a Brownian motion under the tilted measure $Q$) gives
\[
\PP(A_t)\ \ge\ e^{-a^2t/2}\,e^{-a(b_0+1)}\,
Q\bigl(\widetilde W_u\ge b_0\ \forall u\le t,\ \widetilde W_t\in[b_0,b_0+1]
\bigr),
\]
and the $Q$-probability is of order $t^{-3/2}$ by the reflection
principle: polynomial corrections do not affect the rate.
(ii) With $K:=\tfrac{h_0}4t$,
$\PP(T>t)\le e^{-K}+\PP(\Lambda_t\le K)$. On $\{\Lambda_t\le K\}$,
$\Leb\{u\le t:D_u\le D_0\}\le K/(h_0e^{-D_0}\wedge h_0)\lesssim t/4$ (adjust
constants by $D_0$), so there exists
$u_*\in[t/2,\,3t/4+1]$ with $D_{u_*}>D_0$, forcing
$\sigma I_{u_*}\ge|\nu|u_*^2/2-\mu_0u_*\ge|\nu|t^2/16$ for large $t$. By
the Borell--TIS inequality for the Gaussian process $I$ on $[0,t]$ (whose
maximal variance is $t^3/3$ and whose expected supremum is
$O(t^{3/2})=o(t^2)$),
$\PP(\sup_{u\le t}I_u\ge|\nu|t^2/16\sigma)\le
\exp(-c't^4/t^3)=e^{-c't}$; combine.
(iii) Rescale $W_{tv}=\sqrt t\,B_v$, so
$I_{tv}=t^{3/2}\int_0^vB$ and
$D_{tv}=\mu_0tv+\tfrac{\nu t^2v^2}2+\sigma t^{3/2}\!\int_0^vB$. Keeping
$\Lambda_t=O(1)$ requires, at leading order, $D_{tv}\ge0$ for all
$v\in[0,1]$; with $B=\sqrt t\,\gamma$ this reads
$\int_0^v\gamma\ge\beta v^2$, $\beta:=|\nu|/2\sigma$, at
Cameron--Martin cost $\exp(-t\cdot\tfrac12\int_0^1\dot\gamma^2)$. The
variational problem
$\min\tfrac12\int_0^1\dot\gamma^2$ subject to
$\int_0^v\gamma\ge\beta v^2\ \forall v$, $\gamma(0)=0$, is solved by
relaxing to the endpoint constraint
$\int_0^1\gamma=\int_0^1(1-s)\dot\gamma(s)\,ds\ge\beta$: Cauchy--Schwarz
gives $\int\dot\gamma^2\ge3\beta^2$ with equality for
$\dot\gamma(s)=3\beta(1-s)$, i.e.\ $\gamma(v)=3\beta(v-\tfrac{v^2}2)$,
which satisfies the \emph{full} constraint,
$\int_0^v\gamma=\tfrac{\beta v^2}2(3-v)\ge\beta v^2$ on $[0,1]$ with
equality only at $v=1$. Hence the minimum is $\tfrac32\beta^2
=3\nu^2/8\sigma^2$, and unwinding the substitutions yields the profiles
stated: $\mu_{tv}=\mu_0+t(\nu v+\sigma\gamma(v))
\approx\tfrac{|\nu|t}4v(2-3v)$ and
$D_{tv}\approx t^2\bigl(\sigma\int_0^v\gamma-\tfrac{|\nu|v^2}2\bigr)
=\tfrac{|\nu|t^2}4v^2(1-v)$. Elevating this computation to a full
large-deviation proof (upper bound with matching constant, treatment of the
soft constraint) is routine but lengthy and is not carried out here.
\end{proof}

\begin{remark}\label{sh:rem:advdrift}
Three comments. \emph{(1)} The naive lower-bound strategy ``hold
$\mu\ge1$ throughout,'' used in (i), costs $\nu^2/2\sigma^2$ per unit time
and is strictly suboptimal: the true optimizer spends the final third of
the budget with $\mu<0$, spending credit rather than defending the rate ---
survival to a deadline is cheaper than survival in perpetuity, and the
discount is exactly the factor $\tfrac34$. \emph{(2)} Conditioned on the
rare event $\{T>t\}$, the path concentrates on the stated profile: once
again the appearance of purposeful resistance is manufactured by
conditioning. \emph{(3)} The analogue for a mean-reverting driver with
$\bar\mu<0$ is an exponential rate given by the corresponding
Freidlin--Wentzell action; we do not compute it here.
\end{remark}

\begin{center}
\small
\begin{tabular}{@{}lllll@{}}
\toprule
\textbf{Driver} $\mu_t$ & $\PP(\Imm)$ & \textbf{tail of }$\PP(T>t)$ &
$\E[T]$ & \textbf{horizon}\\
\midrule
constant $\mu_0>0$ & $1$ & $\to e^{-h_0/\mu_0}>0$ & $\infty$ (atom) &
trivial step\\
constant $\mu_0<0$ & $0$ & $\exp\bigl(-\tfrac{h_0}{|\mu_0|}
e^{|\mu_0|t}\bigr)$ & finite & ---\\
$\mu_0+\sigma W_t$ & $0$ & $t^{-1/4}$ \textbf{[A]} & $\infty$ & none
(null rec.)\\
$\mu_0+\nu t+\sigma W_t$, $\nu>0$ & $1$ & $\to$ const $>0$ & $\infty$ &
dichotomy in $\nu$\\
$\mu_0+\nu t+\sigma W_t$, $\nu<0$ & $0$ &
$e^{-\frac{3\nu^2}{8\sigma^2}t(1+o(1))}$ \textbf{[S]} & finite & ---\\
OU, mean $\bar\mu>0$ & $1$ & $\to$ const $>0$ & $\infty$ & sharp in
$\bar\mu$\\
OU, mean $\bar\mu=0$ & $0$ & $t^{-1/2}$ \textbf{[S]} & $\infty$ & null
rec., new exponent\\
OU, mean $\bar\mu<0$ & $0$ & $e^{-ct}$ & finite & ---\\
unstable OU \eqref{sh:eq:uou} & $\Phi\bigl(\mu_0\sqrt{2\nu}/\sigma\bigr)$
& mixture & $\infty$ & \textbf{gradual, width }$\sigma/\sqrt{2\nu}$\\
general bounded scale & \eqref{sh:eq:scaleformula} & mixture & $\infty$ &
shape $=$ scale fn.\\
\bottomrule
\end{tabular}
\end{center}

\begin{remark}[Exponents measure smoothness, not recurrence class]
\label{sh:rem:smoothness}
The mean-zero mean-reverting driver is mortal
(Proposition~\ref{sh:prop:recurrent}(c)) with expected tail $t^{-1/2}$, not
$t^{-1/4}$: at large scales its credit is Brownian-like (an additive
functional obeying an invariance principle \cite{Bha82}), so the relevant
persistence exponent is that of Brownian motion, $\tfrac12$, not of its
integral, $\tfrac14$. The survival exponent is a \emph{persistence exponent
of the credit process} and is set by the smoothness of the driver rather
than by its recurrence class: smoother driver, smaller exponent, heavier
survival tail. For a fractional Brownian driver of Hurst index $H$ the
persistence exponent of the driver itself is $1-H$ \cite{Mol99}; the
exponent for its integral --- the relevant one here --- is not known
rigorously \textbf{[O]}: the numerical evidence of Molchan--Khokhlov
\cite{MK04} supports the conjecture $\theta(H)=H(1-H)$ for persistence
from the origin, which correctly returns $\tfrac14$ at $H=\tfrac12$;
see also \cite{AS15}.
\end{remark}

% ---------------------------------------------------------------------
\subsection{The survivorship interpolation and its limit}
\label{sh:sub:interp}

Everything so far has asked where survival is decided in \emph{state
space}; for the Brownian driver the answer was: nowhere. We now change
the question to the one this section exists for. The unbiased law $\PP$
is the ensemble of systems that have survived at least time zero;
conditioning on $\{T>t\}$ produces the ensemble of $t$-survivors; the
object of study is the interpolation between the two and its behavior for
large $t$. The interpolation, it turns out, is where the horizon lives.

\subsubsection{The interpolation is a monotone flow of upward tilts}
\label{sh:subsub:flow}

\begin{definition}[Survivorship interpolation; bias field]
\label{sh:def:interp}
For $t>0$ set $\PP^{(t)}:=\PP(\,\cdot\mid T>t)$, and analogously
$\PP^{(t)}_\tau:=\PP(\,\cdot\mid\tau>t)$ for the hard-wall functional;
$\PP^{(0)}=\PP$ is the unbiased law. For $0\le s<t$ the \emph{bias field}
on the window $[0,s]$ is the Radon--Nikodym derivative
\begin{equation}\label{sh:eq:bias}
\beta_t(s):=\frac{d\PP^{(t)}}{d\PP}\Big|_{\Fc_s}
=\one_{\{T>s\}}\,
\frac{\PP_{X_s}(T>t-s)}{\PP_{X_0}(T>t)},
\qquad X_s=(\mu_s,D_s),
\end{equation}
the second equality by the Markov property; likewise for
$\PP^{(t)}_\tau$ with $\tau$ in place of $T$. The interpolation at finite
$t$ is thus globally absolutely continuous with respect to $\PP$, with
density $\one_{\{T>t\}}/\PP(T>t)$ on $\Fc_\infty$.
\end{definition}

\begin{lemma}[Survivorship tilts upward; \textbf{P}]\label{sh:lem:monotone}
In the driftless model, couple two initial data
$(\mu_0,D_0,h_0)\le(\mu_0',D_0',h_0)$ (componentwise in the first two
arguments) synchronously, with the same $W$. Then pathwise
$\mu'_t-\mu_t=\mu_0'-\mu_0\ge0$,
$D'_t-D_t=(D_0'-D_0)+(\mu_0'-\mu_0)t\ge0$, hence $h'\le h$,
$\Lambda'\le\Lambda$, and both $\PP_x(T>t)$ and $\PP_x(\tau>t)$ are
nondecreasing in $(\mu_0,D_0)$ (and nonincreasing in $h_0$). Consequently
every bias field \eqref{sh:eq:bias} is, on $\{T>s\}$, a nondecreasing
function of the state $X_s$: at every conditioning horizon, survivorship
bias is an upward tilt in rate and credit, and the interpolation is a
monotone flow of such tilts.
\end{lemma}

\begin{proof}
Immediate from the displayed pathwise inequalities and
Definition~\ref{sh:def:cox}.
\end{proof}

\subsubsection{The critical limit is the pure Doob transform}
\label{sh:subsub:Q}

\begin{theorem}[Existence and identification of the limit;
\textbf{C}, soft version \textbf{A}]\label{sh:thm:Q}
Assume the pointwise refined tail asymptotics
$\PP_{(\mu,d)}(\tau>t)\sim c\,\hm(\mu,d)\,t^{-1/4}$ of
Theorem~\ref{sh:thm:persist} and
Proposition~\ref{sh:prop:harmonic}, together with the martingale property
$\E_{x}[\hm(X_s)\one_{\{\tau>s\}}]=\hm(x)$ furnished by \cite{GJW99}.
Then for every $s\ge0$ and every bounded $\Fc_s$-measurable $\Psi$,
\[
\E\bigl[\Psi\mid\tau>t\bigr]\ \longrightarrow\ \E_Q[\Psi],
\qquad t\to\infty,
\]
where $Q$ is the time-homogeneous Markov law with
\begin{equation}\label{sh:eq:Qdensity}
\frac{dQ}{d\PP}\Big|_{\Fc_s}
=\one_{\{\tau>s\}}\,\frac{\hm(X_s)}{\hm(X_0)} :
\end{equation}
the \emph{pure} Doob $\hm$-transform, at eigenvalue one. Under
Ansatz~\ref{sh:ansatz} the same limit holds for
$\PP^{(t)}=\PP(\,\cdot\mid T>t)$.
\end{theorem}

\begin{proof}
Write $f_t:=\one_{\{\tau>s\}}\PP_{X_s}(\tau>t-s)/\PP_{X_0}(\tau>t)$, so
that $\E[\Psi\mid\tau>t]=\E[\Psi f_t]$. By the Markov property
$\E[f_t]=1$ for every $t$. Factorizing,
\[
f_t=\one_{\{\tau>s\}}\,
\frac{(t-s)^{1/4}\,\PP_{X_s}(\tau>t-s)}{t^{1/4}\,\PP_{X_0}(\tau>t)}
\,\Bigl(1-\frac st\Bigr)^{-1/4}
\ \longrightarrow\
f:=\one_{\{\tau>s\}}\frac{\hm(X_s)}{\hm(X_0)}
\quad\text{a.s.},
\]
using the assumed asymptotics at the (fixed, random) starting point
$X_s$ and at $X_0$, and $(1-s/t)^{-1/4}\to1$. The martingale property
gives $\E[f]=1=\lim\E[f_t]$, so by Scheff\'e's lemma $f_t\to f$ in
$L^1(\PP)$, and $\E[\Psi f_t]\to\E[\Psi f]=\E_Q[\Psi]$ for bounded
$\Psi$. That \eqref{sh:eq:Qdensity} defines a consistent, time-homogeneous
Markov family is the standard theory of Doob transforms by invariant
harmonic functions; time-homogeneity with no eigenvalue factor is
possible precisely because $\hm$ is invariant (eigenvalue one) for the
killed semigroup.
\end{proof}

\begin{remark}[Why survivorship has no clock here; \textbf{P}]
\label{sh:rem:noclock}
The mechanism deserves isolation. For \emph{regularly varying} tails the
kinematic factor $\PP_x(\tau>t-s)/\PP_x(\tau>t)\to1$ for each fixed $s$:
the age of the conditioning drops out at leading order, and the limit is
an un-tilted $h$-transform. For \emph{exponential} tails
$\PP_x(T>t)\approx C(x)e^{-\theta t}$ the same factor converges to
$e^{\theta s}$, and survivorship acquires a clock --- the eigen-tilting
of mode (III) in Theorem~\ref{sh:thm:modes}. The polynomial class is
exactly the class in which survivorship bias is, asymptotically, a
change of measure with no time-dependence. In the taxonomy of
penalisations of Wiener measure \cite{RVY06}, the survival weight
$e^{-\Lambda_t}$ thus lands in the clockless (martingale-normalisable)
class precisely because its expectation is regularly varying.
\end{remark}

\begin{proposition}[Interpolation rate; the tail index as linear-response
coefficient; \textbf{P/C}]\label{sh:prop:rate}
Under the hypotheses of Theorem~\ref{sh:thm:Q}, for fixed $s$ as
$t\to\infty$,
\begin{equation}\label{sh:eq:rate-interp}
\beta_t(s)\;=\;\one_{\{\tau>s\}}\,\frac{\hm(X_s)}{\hm(X_0)}
\Bigl(1-\frac st\Bigr)^{-1/4}\bigl(1+o(1)\bigr),
\qquad
\Bigl(1-\frac st\Bigr)^{-1/4}=1+\frac{s}{4t}+O\!\Bigl(\frac{s^2}{t^2}\Bigr).
\end{equation}
The natural interpolation parameter is therefore $s/t$, and the
persistence exponent $\tfrac14$ is the linear-response coefficient of
survivorship bias: the leading finite-$t$ correction to the survivor's
universe is $s/4t$ per window of length $s$. Approach to the limit is
polynomial --- the critical signature, contrasting with the exponential
approach of mode (I) and the exponentially growing tilt of mode (III)
below.
\end{proposition}

\begin{proof}
Insert the assumed asymptotics into \eqref{sh:eq:bias} and expand.
\end{proof}

\subsubsection{The limit object: survivorship drift and promotion to
transience}\label{sh:subsub:promotion}

\begin{theorem}[Survivorship drift; \textbf{P} given \textbf{C} inputs]
\label{sh:thm:drift}
Under $Q$ the state solves
\begin{equation}\label{sh:eq:Qsde}
d\mu_t=\sigma^2\,\partial_\mu\log\hm(\mu_t,D_t)\,dt+\sigma\,dW^Q_t,
\qquad dD_t=\mu_t\,dt,
\end{equation}
for a $Q$-Brownian motion $W^Q$: conditioning adds drift only in the
noisy coordinate. With the representation \eqref{sh:eq:harm},
\begin{equation}\label{sh:eq:G}
b_Q(\mu,D)=\frac{\sigma^2}{\mu}\,G(\kappa),
\qquad
G(\kappa):=\frac12-3\kappa\,\frac{F'(\kappa)}{F(\kappa)},
\end{equation}
and: \emph{(i)} $b_Q\ge0$ everywhere (survivorship never pushes the rate
down), by Lemma~\ref{sh:lem:monotone} applied to $\hm$;
\emph{(ii)} $G(0^+)=\tfrac12$, since
$\hm(\mu,0^+)=c_1'\mu^{1/2}$ by Theorem~\ref{sh:thm:persist};
\emph{(iii)} $G(\kappa)\to0$ as $\kappa\to\infty$, since
$\hm(0^+,d)=c_2 d^{1/6}$ is independent of $\mu$ at leading order.
Statements (ii)--(iii) use the regularity of the profile $F$ furnished by
the explicit formulas of \cite{GJW99}.
\end{theorem}

\begin{proof}
For invariant harmonic $\hm$, the transformed generator is
\[
\Lc^{\hm}f=\hm^{-1}\Lc(\hm f)
=\Lc f+\sigma^2\bigl(\partial_\mu\log\hm\bigr)\,\partial_\mu f,
\]
since only $\partial_\mu^2$ is second order in
\eqref{sh:eq:kolm}; this is \eqref{sh:eq:Qsde}. Formula \eqref{sh:eq:G}
is the logarithmic derivative of \eqref{sh:eq:harm} using
$\partial_\mu\kappa=-3\kappa/\mu$. Positivity: synchronous coupling gives
$\PP_{(\mu,d)}(\tau>t)$ nondecreasing in $\mu$, hence
$\partial_\mu\hm\ge0$ wherever the derivative exists. The boundary values
follow from the quoted prefactor laws: $F(0)=c_1'\in(0,\infty)$ and
$\kappa F'(\kappa)/F(\kappa)\to0$ as $\kappa\to0$ by smoothness of $F$ at
the origin; as $\kappa\to\infty$, $\hm\sim c_2d^{1/6}$ forces
$\mu\,\partial_\mu\log\hm\to0$.
\end{proof}

\begin{corollary}[Wall drift; promotion to transience; \textbf{P/C}]
\label{sh:cor:bessel}
Near zero credit ($\kappa\to0^+$) the survivorship drift reduces to
$\sigma^2/2\mu$, so that in this regime the conditioned rate solves
\[
d\mu_t=\frac{\sigma^2}{2\mu_t}\,dt+\sigma\,dW^Q_t .
\]
Two cautions attach to this formula.
\emph{(i)} This is the Bessel drift of index $\delta=2$, not $\delta=3$:
the exponent $\tfrac12$ in $\hm(\mu,0^+)\propto\mu^{1/2}$ is half the
exponent $1$ of the classical harmonic function $x$ for Brownian motion
killed at the origin, the halving being exactly that of the persistence
exponent from $\tfrac12$ to $\tfrac14$ under one integration.
\emph{(ii)} Under $Q$ the coordinate $\mu$ is not autonomous --- by
\eqref{sh:eq:G} its drift depends on $\kappa$, hence on both coordinates
--- so the comparison is a statement about the \emph{form} of the drift in
a limiting regime, not an identification of processes.

The correct P\'olya statement is therefore not a Bessel index but a
recurrence class, and at that level it is unconditional: under $\PP$ the
credit is null recurrent, returning to its initial level a.s.\
(Theorem~\ref{sh:thm:noimm}) with infinite expected return time
(Corollary~\ref{sh:cor:mean}), whereas under $Q$ it is transient,
$D_s\to\infty$ a.s.\ (Theorem~\ref{sh:thm:singular}). Conditioning on
survival promotes the credit from the recurrent to the transient class ---
the P\'olya $2\mathrm{D}\to3\mathrm{D}$ transition --- with no literal
change of dimension anywhere. Dimension remains, as in
\S\ref{sh:subsub:scale}, the wrong order parameter for the state dynamics;
what survivorship alters is the recurrence class of the \emph{ensemble}.
\end{corollary}

\begin{theorem}[Locally ours, globally elsewhere; \textbf{P/C}]
\label{sh:thm:singular}
\emph{(i)} $Q(\tau=\infty)=1$.
\emph{(ii)} $\hm(X_s)\to\infty$ $Q$-a.s.\ as $s\to\infty$; moreover
$D_s\to\infty$ $Q$-a.s.\ \cite{GJW99}.
\emph{(iii)} $\liminf_sD_s=-\infty$ $\PP$-a.s.
\emph{(iv)} Hence $Q|_{\Fc_s}\ll\PP|_{\Fc_s}$ for every finite $s$, with
the explicit density \eqref{sh:eq:Qdensity}, while $Q\perp\PP$ on
$\Fc_\infty$: the survivor's universe is locally absolutely continuous
with respect to ours on survivors, and globally disjoint from it. The
exit from the measure class is completed only at $t=\infty$; at every
finite conditioning the bias is the bounded tilt of
Definition~\ref{sh:def:interp}.
\end{theorem}

\begin{proof}
(i) $Q(\tau>s)=\E_\PP[\one_{\{\tau>s\}}\hm(X_s)]/\hm(X_0)=1$ for every
$s$ by the martingale property. (ii) Under $Q$, $1/\hm(X_s)$ is a
nonnegative supermartingale:
$\E_Q[1/\hm(X_{s+r})\mid\Fc_s]
=\PP_{X_s}(\tau>r)/\hm(X_s)\le1/\hm(X_s)$; it therefore converges
$Q$-a.s., and since
$\E_Q[1/\hm(X_s)]=\PP_{X_0}(\tau>s)/\hm(X_0)\to0$, Fatou forces the a.s.\
limit to vanish, i.e.\ $\hm(X_s)\to\infty$; the sharper statement
$D_s\to\infty$ is part of the description of the conditioned process in
\cite{GJW99}. (iii) is contained in the proof of
Theorem~\ref{sh:thm:noimm} (the set $A$ there is unbounded, and on it
$D_u\le D_0+\mu_0u-\sigma\varepsilon u^{3/2}\to-\infty$). (iv) Local
absolute continuity is \eqref{sh:eq:Qdensity}; mutual singularity on
$\Fc_\infty$ follows since $\{\lim_sD_s=+\infty\}$ has $Q$-probability
one and $\PP$-probability zero.
\end{proof}

\subsubsection{The three modes of survivorship}\label{sh:subsub:modes}

\begin{theorem}[Modes of the interpolation; \textbf{P/C/S}]
\label{sh:thm:modes}
As $t\to\infty$ the survivorship interpolation exhibits exactly three
asymptotic modes across the driver classes of \S\ref{sh:sub:horizon}:
\begin{enumerate}[label=\emph{(\Roman*)},itemsep=2pt]
\item \emph{Saturation} \textbf{[P]} (favorable drivers,
$\PP(T=\infty)>0$): since $\{T>t\}\downarrow\{T=\infty\}$,
\[
\bigl\|\PP^{(t)}-\PP(\,\cdot\mid T=\infty)\bigr\|_{\mathrm{TV}}
\ \le\ \frac{2\,\PP(t<T<\infty)}{\PP(T>t)}\ \longrightarrow\ 0
\quad\text{on }\Fc_\infty,
\]
and every bias is bounded by $1/\PP(T=\infty)$: survivorship saturates at
a bounded, globally absolutely continuous reweighting. No new universe.
\item \emph{Transformation} \textbf{[C/A]} (critical drivers; the
driftless model): the limit is the pure Doob transform of
Theorem~\ref{sh:thm:Q}, locally absolutely continuous and globally
singular (Theorem~\ref{sh:thm:singular}), approached at the polynomial
rate of Proposition~\ref{sh:prop:rate}. The new universe is entered, but
only at infinity.
\item \emph{Concentration} \textbf{[S]} (adverse drivers, exponential
tails; e.g.\ Proposition~\ref{sh:prop:advdrift}): on fixed windows the
limit is the eigen-tilted transform of
Remark~\ref{sh:rem:noclock} with clock
$\theta=3\nu^2/8\sigma^2$, while at macroscopic scale the conditioned
path law collapses onto the deterministic optimal profile of
Proposition~\ref{sh:prop:advdrift}(iii):
$\sup_{v\in[0,1]}\bigl|\mu_{tv}/t-\tfrac{|\nu|}4v(2-3v)\bigr|\to0$ in
$\PP^{(t)}$-probability. Survivorship becomes a law of large numbers for
the rare: the survivor ensemble is asymptotically deterministic at
leading order.
\end{enumerate}
The three modes are in bijection with the horizon trichotomy
(Theorem~\ref{sh:thm:trichotomy}, Proposition~\ref{sh:prop:recurrent}),
and they constitute a strictly finer invariant: modes (II) and (III)
share $\PP(\Imm)=0$ but are separated by the interpolation --- by whether
the survivor's universe is approached polynomially, without a clock, or
exponentially, with one.
\end{theorem}

\begin{proof}
(I) is continuity of measure from above plus the displayed
total-variation estimate, whose numerator is
$2\PP(\{t<T<\infty\})$ and whose denominator is bounded below by
$\PP(T=\infty)>0$. (II) collects Theorems~\ref{sh:thm:Q}
and~\ref{sh:thm:singular} and Proposition~\ref{sh:prop:rate}. (III) is
the large-deviation sketch of Proposition~\ref{sh:prop:advdrift}(iii)
read under the conditional measure, together with
Remark~\ref{sh:rem:noclock}; elevating it to a full proof requires the
sharp prefactor asymptotics $\PP_x(T>t)\sim C(x)e^{-\theta t}$, which we
do not establish here.
\end{proof}

\begin{remark}[The verdict]\label{sh:rem:verdict}
The section's question can now be answered in one paragraph. The
survivorship interpolation $t\mapsto\PP^{(t)}$ is a monotone flow of
upward tilts (Lemma~\ref{sh:lem:monotone}), globally absolutely
continuous at every finite $t$, whose asymptotic character is the
sharpest invariant of the system and takes exactly three forms
(Theorem~\ref{sh:thm:modes}). The driftless model realizes the critical
form: the limit exists, is the pure Doob transform by the
persistence-harmonic function, carries no clock
(Remark~\ref{sh:rem:noclock}), is approached at rate $s/4t$ with the tail
exponent as linear-response coefficient
(Proposition~\ref{sh:prop:rate}), is locally equivalent on survivors to
the unbiased law yet globally singular
(Theorem~\ref{sh:thm:singular}), and differs in generator by the
universal survivorship drift \eqref{sh:eq:G}, equal to $\sigma^2/2\mu$
at the wall and carrying the credit from the recurrent to the transient
class (Corollary~\ref{sh:cor:bessel}). The gauge structure of
\S\ref{sh:sub:model} persists to the end: the initial data cancel from
the limit density \eqref{sh:eq:Qdensity} up to the normalization
$\hm(X_0)$, and the survivorship drift is universal. The event horizon
absent from state space is therefore real, but it is a feature of the
ensemble, not of the state: it is crossed by the conditioning parameter
$t$ --- which is to say, by survival itself --- and never by the system.
\end{remark}

% ---------------------------------------------------------------------
\subsection{Reference sheet: the filter at a glance}\label{sh:sub:sheet}

This subsection collects, without proof, the quantitative content of the
section in the form most useful for reference. Everything here is
established above; cross-references point to the source.

\subsubsection{The conditioned dynamics in closed form}
\label{sh:subsub:sheet:dyn}

Under the survivorship limit $Q$ of Theorem~\ref{sh:thm:Q} the system is
\begin{equation}\label{sh:eq:sheet}
d\mu_t=\frac{\sigma^2}{\mu_t}\,G(\kappa_t)\,dt+\sigma\,dW^Q_t,
\qquad
dD_t=\mu_t\,dt,
\qquad
\kappa=\frac{\sigma^2D}{\mu^3},
\end{equation}
with $G(\kappa)=\tfrac12-3\kappa F'(\kappa)/F(\kappa)$ read off
$\hm=\mu^{1/2}F(\kappa)$ as in \eqref{sh:eq:harm}--\eqref{sh:eq:G}.
Dimensionally $[\sigma^2/\mu]=\mathsf T^{-2}=[\mu]/[t]$, as required. The
shape function interpolates between two endpoints
(Theorem~\ref{sh:thm:drift}):
\[
G(0^+)=\tfrac12\ \ (\text{drift }\sigma^2/2\mu,\ \text{maximal}),
\qquad
G(\infty)=0\ \ (\text{drift }0,\ \text{none}),
\]
so that \emph{survivorship pushes hardest near the wall and not at all
once credit is abundant}, the strength along any ray of fixed $\kappa$
scaling as $1/\mu$.

\begin{remark}[The endpoints are ray limits, not coordinate limits]
\label{sh:rem:raylimit}
One may not simply let $\mu\to0$ in $\sigma^2/2\mu$: at fixed $D>0$ that
sends $\kappa\to\infty$ and hence $G\to0$, and the two factors compete.
Only the statement along rays of constant $\kappa$ is unambiguous, which
is precisely why $\kappa$ and not $(\mu,D)$ separately is the coordinate
in which the conditioned dynamics should be read
(Proposition~\ref{sh:prop:rays}).
\end{remark}

\subsubsection{The three levels, before and after}
\label{sh:subsub:sheet:levels}

\begin{center}\small
\begin{tabular}{@{}p{0.24\textwidth}p{0.34\textwidth}p{0.34\textwidth}@{}}
\toprule
& \textbf{Before} ($\PP$) & \textbf{After} ($Q$)\\
\midrule
\multicolumn{3}{@{}l}{\emph{Level 1: the driver $\mu$}}\\
law & Gaussian, driftless & non-Gaussian, drift
$(\sigma^2/\mu)G(\kappa)\ge0$\\
autonomy & $\mu$ Markov by itself & only the pair $(\mu,D)$ is Markov\\
quadratic variation & $\sigma^2\,dt$ & $\sigma^2\,dt$ \emph{(unchanged)}\\
increments & stationary, independent & neither\\
self-similarity & index $3/2$ & index $3/2$ \emph{(preserved)}\\
\addlinespace
\multicolumn{3}{@{}l}{\emph{Level 2: the hazard $h$}}\\
path regularity & $C^1$ & $C^1$ \emph{(unchanged)}\\
law at time $s$ & lognormal, $\log$-variance
$\sigma^2s^3/3$ & non-Gaussian, one-sided\\
range & exceeds $h_0$ infinitely often,
unbounded above & $h_s<h_0$ for \emph{all} $s$ (if $D_0=0$)\\
credit $D$ & null recurrent,
$\liminf_sD_s=-\infty$ & transient, $D_s\to\infty$,
$D_s\asymp\sigma s^{3/2}$\\
\addlinespace
\multicolumn{3}{@{}l}{\emph{Level 3: survival}}\\
$\Lambda_\infty$ & $=\infty$ a.s. & $<\infty$ a.s.,
$\asymp\hat h_0=h_0\sigma^{-2/3}$\\
quenched curve $e^{-\Lambda_t}$ & $\to0$; collapse
$\exp(-e^{\sigma t^{3/2}})$ & plateau at $e^{-\Lambda_\infty}>0$\\
population curve & $\asymp Ct^{-1/4}$ & $\equiv1$\\
lifetime & $T<\infty$ a.s., $\E[T]=\infty$ & $T=\infty$\\
\bottomrule
\end{tabular}
\end{center}

\noindent Sources: Level 1, Theorems~\ref{sh:thm:Q}
and~\ref{sh:thm:drift}; Level 2, Definition~\ref{sh:def:model} and
Theorem~\ref{sh:thm:singular}; Level 3, Theorem~\ref{sh:thm:noimm},
Claim~\ref{sh:claim:tail}, Corollary~\ref{sh:cor:mean}.

\begin{remark}[Notable structures at Level 2]\label{sh:rem:sheet:L2}
Three features are worth isolating. \emph{(i)} The hazard is one
derivative \emph{smoother} than the noise driving it, a direct
consequence of the log-derivative placement
(Remark~\ref{sh:rem:iwp}); conditioning does not change this, since
Girsanov alters drift and never path regularity. \emph{(ii)} Before
conditioning $\log h_s$ is exactly Gaussian with variance
$\sigma^2s^3/3$ --- \emph{cubic} in time, the integrated-Brownian
fingerprint --- so $h_s$ is lognormal. \emph{(iii)} After conditioning,
the survivor's hazard never once exceeds its initial value, whereas
before it does so infinitely often and by unbounded amounts: this is the
wall of \S\ref{sh:subsub:reduction} in its most concrete form. Finer
still, Groeneboom, Jongbloed and Wellner \cite{GJW99} identify
$t\mapsto t^{9/10}$ as a critical curve for the conditioned process, the
expected time spent below $t^\alpha$ being finite for $\alpha<9/10$ and
infinite for $\alpha\ge9/10$: survivors do lag their typical
$s^{3/2}$ growth, and $9/10$ is exactly how much.
\end{remark}

\begin{proposition}[The quenched--annealed gap; \textbf{P}]
\label{sh:prop:gap}
In the driftless model, almost surely
\begin{equation}\label{sh:eq:gap}
\log\Lambda_t=\sup_{u\le t}\bigl(D_0-D_u\bigr)+O(\log t),
\end{equation}
and $t^{-3/2}\sup_{u\le t}(D_0-D_u)\to\sigma\sup_{v\le1}(-I_v)$ in law, a
nondegenerate positive limit. Hence the quenched survival curve
collapses \emph{doubly exponentially}: writing
$c_t:=t^{-3/2}\log\Lambda_t$,
\[
e^{-\Lambda_t}=\exp\bigl(-e^{\,c_t\,t^{3/2}}\bigr),
\qquad
c_t\ \xrightarrow{\;d\;}\ \sigma\sup_{v\le1}(-I_v)\;>\;0,
\]
while by Claim~\ref{sh:claim:tail} the population curve decays only as
$Ct^{-1/4}$. The entire population tail is therefore carried by a set of
paths of vanishing probability.
\end{proposition}

\begin{proof}
Write $M_t=\sup_{u\le t}(D_0-D_u)$ and
$\Lambda_t=h_0\int_0^te^{D_0-D_u}du$. The upper bound
$\Lambda_t\le h_0te^{M_t}$ is immediate. For the lower bound, $D$ is
$C^1$ with $\dot D=\mu$ and $\sup_{u\le t}|\mu_u|=O(t^{1/2})$ a.s.\ up to
logarithmic factors, so on an interval of length $\asymp t^{-1/2}$ around
a near-maximiser the integrand exceeds $e^{M_t-1}$, giving
$\Lambda_t\gtrsim h_0t^{-1/2}e^{M_t}$. Taking logarithms yields
\eqref{sh:eq:gap}. The scaling limit is
Proposition~\ref{sh:prop:selfsim} applied to $I$, with
$\sup_{v\le1}(-I_v)>0$ a.s.
\end{proof}

\begin{remark}[The interpolation as a one-parameter family]
\label{sh:rem:sheet:family}
Between the two columns, the survival curve of a $t$-survivor is
\[
\PP(T>t+s\mid T>t)\;\approx\;\Bigl(1+\frac st\Bigr)^{-1/4},
\]
a scale-free family flattening from a power law at $t=0$ to the constant
$1$ as $t\to\infty$; equivalently the $t$-survivor carries effective
hazard $1/4t$ (Corollaries~\ref{sh:cor:residual}
and~\ref{sh:cor:lindy}). The filter does not deform the survival curve
arbitrarily: it slides it along a single one-parameter orbit.
\end{remark}

\subsubsection{What conditioning can and cannot do}
\label{sh:subsub:sheet:rigid}

\begin{proposition}[Rigidity of the survivorship transform; \textbf{P}]
\label{sh:prop:rigid}
Let $\widetilde\PP$ be any measure locally absolutely continuous with
respect to $\PP$ on each $\Fc_s$ --- in particular $Q$. Then:
\begin{enumerate}[label=\emph{(\roman*)},itemsep=1pt]
\item \emph{The kinematics survive exactly.} $dD_t=\mu_t\,dt$ holds
$\widetilde\PP$-a.s.\ unchanged: credit remains literally the integral of
rate. Conditioning cannot alter it.
\item \emph{Only the driver may be tilted.} Since the noise enters solely
the $\mu$ coordinate in \eqref{sh:eq:kolm}, Girsanov can add drift only
to $\mu$; the transform has exactly one functional degree of freedom.
\item \emph{Self-similarity fixes its form.} If in addition
$\widetilde\PP$ respects the scaling of
Proposition~\ref{sh:prop:selfsim}, that degree of freedom is forced into
the shape $(\sigma^2/\mu)G(\kappa)$ of \eqref{sh:eq:sheet}, a single
function of one dimensionless variable.
\item \emph{The cascade is destroyed.} Under $\PP$ the system is
triangular: $\mu$ is autonomous and $D$ is slaved to it. Under $Q$ the
credit re-enters the drift of $\mu$ through $\kappa$, so $\mu$ alone is
no longer Markov although the pair remains so.
\end{enumerate}
\end{proposition}

\begin{proof}
(i) and (ii) are Girsanov's theorem for \eqref{sh:eq:kolm}, whose
diffusion matrix is degenerate with range spanned by $\partial_\mu$; the
finite-variation coordinate admits no equivalent change. (iii) is
Proposition~\ref{sh:prop:harmonic} together with
Theorem~\ref{sh:thm:drift}. (iv) is the $\kappa$-dependence in
\eqref{sh:eq:G}, non-constant by Theorem~\ref{sh:thm:drift}(ii)--(iii).
\end{proof}

\begin{remark}[Survivorship coupling is collider bias in continuous time]
\label{sh:rem:collider}
The feedback of Proposition~\ref{sh:prop:rigid}(iv) is informational, not
mechanical: nothing pushes $\mu$. Under $\PP$ the credit accumulated by
time $s$ and the future increments of the driver are independent, but
both feed a common outcome, survival; conditioning on that outcome
renders them dependent. This is exactly collider bias --- Berkson's
paradox, or ``explaining away'' --- realised in continuous time, and
\eqref{sh:eq:sheet} is its quantitative form: little credit banked
implies the remaining path must work harder, and the drift says by how
much. The same object appears in large-deviation theory as the driven or
effective process associated with a rare event. The moral of
Proposition~\ref{sh:prop:rigid} is then sharp: survivorship cannot bend
the kinematics of a system, only the statistics of what drives it ---
and, when the system is self-similar, only along a one-parameter family
of shapes.
\end{remark}

% ---------------------------------------------------------------------
\subsection{Numerical verification}\label{sh:sub:numerics}

The one load-bearing heuristic, Ansatz~\ref{sh:ansatz}, admits a direct
test, which we specify for completeness; this subsection may be read as an
appendix.

\begin{lemma}[Exact simulation of the pair; \textbf{P}]\label{sh:lem:exact}
For a step $\Delta>0$ let $\xi:=W_{t+\Delta}-W_t$ and
$\eta:=\int_t^{t+\Delta}(W_s-W_t)\,ds$. Then $(\xi,\eta)$ is independent of
$\Fc_t$, centered Gaussian with covariance
$\bigl(\begin{smallmatrix}\Delta&\Delta^2/2\\
\Delta^2/2&\Delta^3/3\end{smallmatrix}\bigr)$, and the update
\[
\mu_{t+\Delta}=\mu_t+\sigma\xi,\qquad
D_{t+\Delta}=D_t+\mu_t\Delta+\sigma\eta
\]
reproduces the law of \eqref{sh:eq:kolm} at the grid points with zero
discretization error.
\end{lemma}

\begin{proof}
$(\xi,\eta)$ is a measurable functional of the post-$t$ increments and is a
copy of $(W_\Delta,I_\Delta)$; apply Lemma~\ref{sh:lem:secondorder}. The
update identity is
$I_{t+\Delta}=I_t+W_t\Delta+\eta$ multiplied through by $\sigma$.
\end{proof}

The remaining error is the quadrature of
$\Lambda_t=h_0e^{-D_0}\int_0^te^{-(D_u-D_0)}du$, whose integrand varies on
the scale $|\Delta D|\lesssim1$, i.e.\ $\Delta\lesssim1/|\mu|\sim
t^{-1/2}$: the constraint tightens with time. Accumulate in logarithmic
space; the integrand ranges over hundreds of orders of magnitude. To reach
the tail --- probabilities of $10^{-3}$ to $10^{-4}$ require
$t/t_*$ up to $10^{12}$ and beyond --- use multilevel splitting along the
thresholds $D_k=\sigma(\theta^kt_*)^{3/2}$, or importance sampling tilted
by the harmonic function of Proposition~\ref{sh:prop:harmonic}, the
asymptotically optimal tilt. Validate the code first on the hard-wall
functional, where \eqref{sh:eq:persist} is a theorem; then run the
soft-killing functional and fit $\log\PP$ against $\log t$ over at least
three decades, discarding $t\lesssim10\,t_*$.
\emph{Falsifiers:} a fitted slope robustly different from $-\tfrac14$
(Ansatz fails: the killing term carries its own scale); any dependence of
the slope on $h_0$ (gauge violation, contradicting
Proposition~\ref{sh:prop:gauge}); non-convergence of the conditioned
rescaled credit $D_s/\sigma s^{3/2}$ (failure of the self-similar Yaglom
picture of Remark~\ref{sh:rem:doob}). The survivorship results of
\S\ref{sh:sub:interp} add three sharper tests: estimate the window bias
$\beta_t(s)$ under splitting and regress it against $\hm(X_s)$, whose
proportionality tests Theorem~\ref{sh:thm:Q}; regress conditioned
increments $\Delta\mu$ on $1/\mu$ at small $\kappa$, whose slope
$\sigma^2/2$ tests the wall asymptotics of
Corollary~\ref{sh:cor:bessel}; and verify the linear response
$\beta_t\approx1+s/4t$ at large $t/s$, whose coefficient $\tfrac14$ tests
Proposition~\ref{sh:prop:rate}.

% ---------------------------------------------------------------------
\subsection*{Section summary}

The system is a Kolmogorov diffusion in $(\mu,D)$ with exponential
killing; at exponent level it is the persistence problem for integrated
Brownian motion. Initial data are gauge: they act only through an
equivalent timescale $t_*$, and the one intrinsic coordinate is
$\kappa=\sigma^2D/\mu^3$ (\S\ref{sh:sub:model}). For the driftless model,
viability is impossible for every initial condition
(Theorem~\ref{sh:thm:noimm}), yet life expectancy is infinite:
$\PP(T>t)\asymp(t_*/t)^{1/4}$, with Pareto($\tfrac14$) residual life,
quantiles growing as the fourth power of rarity, and an effective Lindy
hazard $1/4t$ --- null recurrence, the two-dimensional case of the P\'olya
analogy (\S\ref{sh:sub:verdict}). No horizon exists there: the crossover
is a maximally soft quarter-power law smeared over orders of magnitude,
and recovery cost is sublinear in credit. Where genuine horizons exist,
the order parameter is the boundedness of the driver's scale function,
the horizon profile \emph{is} the normalized scale function
(Theorem~\ref{sh:thm:trichotomy}), the recurrent class is decided by the
stationary mean with mortal critical case
(Proposition~\ref{sh:prop:recurrent}), and the solvable unstable-OU
driver realizes the profile exactly as a Gaussian sigmoid
$\Phi(\mu_0\sqrt{2\nu}/\sigma)$ of width $\sigma/\sqrt{2\nu}$, a step
function only in the small-noise limit (\S\ref{sh:sub:horizon}). Under
adverse drift the survival tail is exponential with rate
$3\nu^2/8\sigma^2$, carried by a front-load-then-coast optimal
fluctuation. The destination is the survivorship interpolation
(\S\ref{sh:sub:interp}): a monotone flow of upward tilts $\PP^{(t)}$
whose critical limit is the pure Doob transform by the
persistence-harmonic function, approached without a clock at rate
$s/4t$ --- the tail exponent as linear-response coefficient --- locally
absolutely continuous on survivors yet globally singular, and generated
by the universal survivorship drift, equal to $\sigma^2/2\mu$ at the
wall: conditioning carries the credit from the recurrent to the transient
class, the P\'olya promotion effected by survival rather than by
dimension. Across
drivers the interpolation has exactly three modes --- saturation,
transformation, concentration --- in bijection with the horizon
trichotomy and strictly finer than the survival dichotomy. A horizon is
not a place. Where it exists in state space it is a ratio of scales; and
where it does not, it is real nonetheless --- located in the ensemble,
crossed not by the state but by the conditioning parameter, which is to
say by survival itself.

\section{The two-channel muffled hazard}\label{tc:sec}

\begin{quote}\small
Section~\ref{sh:sec} studied a hazard whose negative logarithmic derivative is
Brownian, and found the event horizon to be soft in state space and real
only in the ensemble. That model was a simplification of a richer
ontology, supplied here. The hazard is not primitive: it splits into two
\emph{channels}. A bounded \emph{crest} channel
$B_t\bigl(1-\tanh\tfrac{x_t}{2}\bigr)$ carries a slowly varying
background $B$ and can shield from it but never amplify it; a
\emph{trough} channel $\Gamma e^{-\rho x_t}$ with brittleness $\rho>1$
supersedes the background entirely at negative $x$ and carries internal
damage. Four results organise the section. First, the background, built
as a sum of exponentiated Brownian motions to force positivity, is a
partition function whose temperature falls with age: it freezes at
$t_{\mathrm{fr}}=2\log n/\epsilon^{2}$, is Gompertz-exponential before
and extreme-dominated after, and is absorbed \emph{exactly} into Section
One's initial data as a constant headwind on the improvement rate
(\S\ref{tc:sub:background}). Second, strength and anatomy are not
parameters but \emph{coordinates}: the state $(\mu_t,x_t)$ transforms
exactly and adaptedly to polar form $x_t=\Ac_t\cos\Theta_t$ with
amplitude $\Ac_t=\sigma t^{3/2}R_t$, where $R_t$ is Rayleigh, $\Theta_t$
uniform, and the two independent \emph{at each fixed time} --- though,
as we show, not as processes (\S\ref{tc:sub:polar}). Third, between the
channels' walls lies a \emph{waiting room} within which the hazard is
the background alone and the system's driver is muted; the model leaves
it at $t_{\mathrm{act}}\asymp(k/\rho\sigma)^{2/3}$ of its own accord, so
no capacity parameter is required (\S\ref{tc:sub:room}). Fourth, with
constant background the whole construction is Section~\ref{sh:sec}'s model with
the credit passed through a soft kink, revealing that the original
driver was the new driver times a state-dependent gain --- one that
vanishes in the void (\S\ref{tc:sub:gain}). Both walls are $o(t^{3/2})$
and hence gauge, so the persistence exponent $\tfrac14$, the infinite
life expectancy, and the survivorship interpolation survive exactly:
Section~\ref{sh:sec} is undisturbed as a theorem, not as an approximation
(\S\ref{tc:sub:asymptotics}).
\end{quote}

\subsection*{Overview and status of results}

Tags are as in Section~\ref{sh:sec}: \textbf{[P]} proved here; \textbf{[C]}
classical or cited; \textbf{[A]} conditional on moving-boundary
robustness of the same class as that section's Ansatz; \textbf{[S]}
scaling-level; \textbf{[O]} open. The chain is: the two-channel model
reduces to Section~\ref{sh:sec} on an explicit slice \textbf{[P]}; the crest
channel is bounded, enforcing the causal asymmetry structurally
\textbf{[P]}; the background has a freezing transition \textbf{[C]}
whose delocalised phase is Gompertz \textbf{[P]}; a linear background is
absorbed exactly into initial data \textbf{[P]} and a sublinear one is
invisible \textbf{[A]}; the polar decomposition is exact and adapted
with Rayleigh/uniform independent marginals \textbf{[P]}, but the polar
\emph{processes} are dependent \textbf{[P]}; bad anatomy occupies
exactly half of logarithmic time \textbf{[P]}; the gain is exact with
slope range $(0,\rho)$ \textbf{[P]}; and Section~\ref{sh:sec}'s asymptotics are
preserved \textbf{[A]}.

\begin{convention}\label{tc:conv}
Notation and the probabilistic setting are as in Section~\ref{sh:sec}, with one
change forced by the reassignment of names: the survival function is
written $\bar F_t=e^{-\Lambda_t}$, freeing $\Sigma$ and $S$ for
strength. $W$ drives $\mu$; $W^{(1)},\dots,W^{(n)}$ are independent
standard Brownian motions independent of $W$ driving the background.
References to Section~\ref{sh:sec} are textual and become cross-references on
merge.
\end{convention}

% ---------------------------------------------------------------------
\subsection{The two-channel model}\label{tc:sub:model}

\subsubsection{Primitives and the hazard}\label{tc:subsub:primitives}

\begin{definition}[Two-channel model]\label{tc:def:model}
Fix $\sigma>0$, $\epsilon>0$, $n\in\mathbb N$, $\Gamma>0$ and
$\rho>1$. Define
\begin{align}
\text{driver:}&\quad \mu_t=\mu_0+\sigma W_t,
&&[\mu]=\mathsf T^{-1},\ [\sigma]=\mathsf T^{-3/2},
\label{tc:eq:driver}\\
\text{muffling exponent:}&\quad x_t=x_0+\int_0^t\mu_s\,ds,
&&[x]=1,
\label{tc:eq:x}\\
\text{background:}&\quad B_t=\sum_{i=1}^{n}e^{\epsilon W^{(i)}_t},
&&[B]=\mathsf T^{-1},\ [\epsilon]=\mathsf T^{-1/2},
\label{tc:eq:B}
\end{align}
and let the total hazard be the sum of two channels,
\begin{equation}\label{tc:eq:hazard}
h^{\mathrm{tot}}_t
=\underbrace{B_t\,\varsigma_b(x_t)}_{\text{crest}}
+\underbrace{\Gamma\,e^{-\rho x_t}}_{\text{trough}},
\qquad
\varsigma_b(x):=\frac{1+e^{-b}}{1+e^{x-b}}
=\frac{1+e^{-b}}{2}\Bigl(1-\tanh\tfrac{x-b}{2}\Bigr),
\end{equation}
with \emph{offset} $b>0$. The normalising factor $1+e^{-b}$ is chosen so
that $\varsigma_b(0)=1$ exactly, whence $h^{\mathrm{crest}}|_{x=0}=B_t$:
the void limit is the background itself, for every $b$. Setting $b=0$
recovers the unshifted form $B_t(1-\tanh\frac{x}{2})$
Survival is $\bar F_t=e^{-\Lambda_t}$ with
$\Lambda_t=\int_0^t h^{\mathrm{tot}}_u\,du$, and the death time $T$ is
given by the Cox construction of Section~\ref{sh:sec}. There is \emph{no}
capacity parameter and \emph{no} leak: $x$ is the undamped integral of
$\mu$, exactly as in Section~\ref{sh:sec}.
\end{definition}

\begin{remark}[Regimes]\label{tc:rem:regimes}
The three regimes of \eqref{tc:eq:hazard}:
\begin{center}\small
\begin{tabular}{@{}lll@{}}
\toprule
regime & total hazard & reading\\
\midrule
$x\gg b$ & $\approx B_te^{\,b-x}$ & shielding; background passes through
multiplicatively\\
$-k/\rho\ll x\ll b$ & $\approx B_t$ & \emph{void}: the system
contributes nothing\\
$x\ll0$ & $\approx\Gamma e^{\rho|x|}$ & internal damage; $B$ has dropped
out entirely\\
\bottomrule
\end{tabular}
\end{center}
The middle line is exact at $x=0$: the normalisation gives
$h^{\mathrm{crest}}|_{x=0}=B_t$, so the void limit is the background
itself. The offset $b$ places $x=0$ \emph{inside} the void rather than
at its upper edge --- see Remark~\ref{tc:rem:whyoffset}.
\end{remark}

\subsubsection{The four constraints}\label{tc:subsub:constraints}

The model is answerable to four structural requirements. We state them
as the author posed them and record where each is met.

\begin{constraint}[Background reaches the total at positive $x$]
\label{tc:C1}
Met for $\rho>1$, and \emph{only} then: for $x\gg b$ the crest channel
is $B_te^{\,b-x}$, multiplicative in $B$, and by
Proposition~\ref{tc:prop:dropout}(ii) it dominates the trough at all
large positive $x$ precisely when $\rho>1$. This --- not
Constraint~\ref{tc:C2} --- is what the brittleness condition buys.
\end{constraint}

\begin{constraint}[Background drops out at negative $x$]\label{tc:C2}
Met by Proposition~\ref{tc:prop:dropout}(i), for \emph{every} $\rho>0$.
Since $\Gamma$ is a primitive constant the trough contains no $B$
whatsoever, so at $x\ll0$ the background disappears from the total
entirely, not merely in its fluctuations.
\end{constraint}

\begin{constraint}[Section~\ref{sh:sec} undisturbed]\label{tc:C3}
Met as a theorem, \S\ref{tc:sub:asymptotics}: both walls are
$o(t^{3/2})$ and hence gauge.
\end{constraint}

\begin{constraint}[No smuggled primitives]\label{tc:C4}
Met by Theorem~\ref{tc:thm:gain}: with $B$ constant the total hazard is
$h^{\mathrm{tot}}(0)e^{-\Psi(x)}$ for a single soft-kink gain $\Psi$, so
the model is Section~\ref{sh:sec}'s with the credit passed through a kink.
\end{constraint}

\begin{proposition}[Which channel dominates, and where; \textbf{P}]
\label{tc:prop:dropout}
\emph{(i) Negative $x$.} For every $\rho>0$,
$\Gamma e^{-\rho x}/\bigl(B\varsigma_b(x)\bigr)\to\infty$ as
$x\to-\infty$, since $\varsigma_b$ is bounded while $e^{\rho|x|}$ is
not. Hence $h^{\mathrm{tot}}\sim\Gamma e^{\rho|x|}$, free of $B$:
Constraint~\ref{tc:C2} needs no condition on $\rho$ beyond positivity.

\emph{(ii) Positive $x$.} As $x\to+\infty$ the crest is
$\sim Be^{\,b-x}$ and the trough is $\Gamma e^{-\rho x}$, so the crest
dominates for all large $x$ if and only if $\rho>1$. For $\rho<1$ the
two cross at
\begin{equation}\label{tc:eq:crossover}
x_{\times}=\frac{k+b}{1-\rho},\qquad k:=\log\frac{B}{\Gamma},
\end{equation}
beyond which the trough dominates and the background drops out
\emph{even for survivors}.
\end{proposition}

\begin{proof}
(i) $\varsigma_b\le1+e^{-b}$ by Proposition~\ref{tc:prop:bounded}, while
$\Gamma e^{\rho|x|}\to\infty$ for any $\rho>0$. (ii) Compare exponents
$b-x$ and $-k-\rho x$; they are equal at \eqref{tc:eq:crossover}, and
for $\rho>1$ the crest exponent is the larger for all $x$ exceeding a
finite threshold.
\end{proof}

\begin{remark}[What $\rho=1$ is, and is not]\label{tc:rem:rhorole}
Three separate claims must be kept apart, and an earlier draft of this
section conflated the first two.
\emph{(a)} $\rho=1$ is \emph{not} the threshold for the background to
drop out at negative $x$; that happens for every $\rho>0$, by
(i).
\emph{(b)} $\rho=1$ \emph{is} the threshold at which the background
ceases to reach the survivor's hazard: for $\rho<1$ internal damage
overtakes the external channel beyond $x_{\times}$ and the system is
effectively closed to the external world, violating
Constraint~\ref{tc:C1}. This is a sharp threshold in a structural
parameter, and it concerns \emph{openness}, not survival.
\emph{(c)} $\rho$ does \emph{not} affect the survival exponent. For any
$\rho>0$ the dominant channel is an exponential in $x$, so the killing
wall is a level of $x$ and the persistence problem is Section~\ref{sh:sec}'s up
to a rescaling of the credit by a constant. The exponent is $\tfrac14$
throughout, which strengthens Constraint~\ref{tc:C3} rather than
qualifying it.
\end{remark}

\subsubsection{Boundedness and the causal asymmetry}
\label{tc:subsub:asymmetry}

\begin{proposition}[The crest channel is bounded; \textbf{P}]
\label{tc:prop:bounded}
$0<\varsigma_b(x)<1+e^{-b}$ for all $x\in\R$, the supremum being
approached only as $x\to-\infty$. Hence the external channel can
\emph{reduce} the background hazard without bound, but can
\emph{increase} it by at most the factor $1+e^{-b}$, which tends to $1$
as $b$ grows.
\end{proposition}

\begin{proof}
$(1+e^{-b})/(1+e^{x-b})$ is decreasing in $x$ with limits $1+e^{-b}$ and
$0$.
\end{proof}

\begin{remark}[Why this matters]\label{tc:rem:asymmetry}
The slogan is the author's: \emph{strength can shield from the
background but never amplify it}, and it is now enforced structurally
rather than assumed. The earlier multiplicative form $B_te^{-x}$
violated it --- at $x<0$ it multiplied the background without bound,
which asserts that a failure of the system's own muffling magnifies the
external world, a causal claim the ontology does not license.
Amplification now arises only through the trough channel, which is
internal damage: a distinct mechanism, correctly separated from the
external one. Proposition~\ref{tc:prop:dropout} is the formal expression
of that separation. Note how much the offset strengthens the slogan: at
$b=0$ the background could still be doubled, whereas for $b\gtrsim6$ the
permitted amplification is under a quarter of one percent. Shielding is
unbounded; amplification is, for practical purposes, forbidden.
\end{remark}

\subsubsection{Reduction to Section~\ref{sh:sec}}\label{tc:subsub:reduction}

\begin{proposition}[Reduction; \textbf{P}]\label{tc:prop:reduction}
On the slice $\epsilon=0$ (so $B\equiv n$), $\Gamma\to0$, and $x\gg0$,
the model is Section~\ref{sh:sec}'s with $h_0=2n$ and credit $x-x_0$. More
usefully: for $\epsilon=0$ and any $\Gamma$, the total hazard is a
deterministic function of $x$ alone, and Theorem~\ref{tc:thm:gain}
identifies it as Section~\ref{sh:sec}'s hazard composed with a fixed gain.
\end{proposition}

\begin{proof}
At $\epsilon=0$, $B\equiv n$; for $x\gg0$, $h^{\mathrm{tot}}\approx
2ne^{-x}$ since the trough is $O(e^{-\rho x})=o(e^{-x})$.
\end{proof}

% ---------------------------------------------------------------------
\subsection{The background as a free energy}\label{tc:sub:background}

The background must be positive, variable, and never near zero.
Construction \eqref{tc:eq:B} achieves all three: a sum of $n$ positive
terms is small only if \emph{all} are small, an event of probability
$\Phi(-a/\epsilon\sqrt t)^{n}$, exponentially small in $n$. This is the
mechanism the author identified. What was not anticipated is that the
construction is a familiar object.

\subsubsection{Freezing and Gompertz}\label{tc:subsub:freezing}

\begin{remark}[The background is a partition function]
\label{tc:rem:partition}
Writing $\epsilon W^{(i)}_t=\epsilon\sqrt t\,g_i$ with $g_i$ i.i.d.\
standard normal,
\[
B_t=\sum_{i=1}^{n}e^{\beta_t g_i},\qquad \beta_t:=\epsilon\sqrt t,
\]
the partition function of $n$ i.i.d.\ Gaussian energy levels at inverse
temperature $\beta_t$. Since $\beta_t$ \emph{increases} with $t$, the
background is a disordered system that \emph{cools as it ages}: it is
Derrida's random energy model \cite{Der81,Der80} with time as inverse
temperature. The corresponding limit theory for sums of random
exponentials is \cite{BBM05}; the mathematical treatment of the REM is
\cite{Bov06}.
\end{remark}

\begin{theorem}[Freezing; \textbf{C}]\label{tc:thm:freezing}
With $\theta=\epsilon\sqrt{t/\log n}$, as $n\to\infty$ at fixed
$\theta$,
\begin{equation}\label{tc:eq:freezing}
\frac{\log B_t}{\log n}\;\longrightarrow\;
\begin{cases}
1+\theta^{2}/2, & \theta<\sqrt2\quad(\text{delocalised}),\\
\theta\sqrt2, & \theta\ge\sqrt2\quad(\text{frozen}),
\end{cases}
\end{equation}
in probability, with freezing time
$t_{\mathrm{fr}}=2\log n/\epsilon^{2}$.
\end{theorem}

\begin{proof}[Proof and status]
The two candidates are the annealed free energy
$\log\E B_t=\log n+\epsilon^{2}t/2$ and the extremal contribution
$\max_i\epsilon W^{(i)}_t\approx\epsilon\sqrt{2t\log n}$; normalised
these are $1+\theta^{2}/2$ and $\theta\sqrt2$, and
$(1+\theta^{2}/2)-\theta\sqrt2=\tfrac12(\theta-\sqrt2)^{2}\ge0$ with
equality only at $\theta=\sqrt2$ --- the annealed value dominates
everywhere and the curves meet tangentially, the REM's characteristic
contact. The second moment gives
$\E B_t^{2}/(\E B_t)^{2}=e^{\epsilon^{2}t}/n+(n-1)/n\to1$ iff
$\epsilon^{2}t<\log n$, so concentration is immediate for
$t<t_{\mathrm{fr}}/2$; the full delocalised range needs the truncated
second moment, and the frozen phase is the standard extreme-value
computation. Both are classical \cite{Der81,Bov06,BBM05}.
\end{proof}

\begin{remark}[The author's two mechanisms are the two phases]
\label{tc:rem:twophases}
The brief describes the background as held up both by ``the high ones
dominating'' and by ``the fuzz of all the intermediate-sized ones adding
together.'' These are not two descriptions of one regime but exact
descriptions of the two: the frozen phase is carried by the single
largest term, the delocalised phase by the bulk.
\end{remark}

\begin{corollary}[Gompertz; \textbf{P}]\label{tc:cor:gompertz}
For $t<t_{\mathrm{fr}}$, $B_t=n\,e^{\epsilon^{2}t/2}(1+o(1))$ in
probability: the background grows exponentially in age at rate
$\epsilon^{2}/2$. This is the Gompertz law \cite{Gom25}, arriving
unbidden --- the construction was chosen for positivity, not to
reproduce it.
\end{corollary}

\subsubsection{The background is gauge}\label{tc:subsub:gauge}

\begin{lemma}[Channels define moving walls; \textbf{P}]
\label{tc:lem:walls}
Set $D_t:=x_t-x_0$. Each channel becomes $O(1)$ --- i.e.\ contributes
appreciably to $\Lambda$ --- at a \emph{killing wall}:
\begin{equation}\label{tc:eq:killwalls}
x^{\mathrm{crest}}_t=\log B_t+b,
\qquad
x^{\mathrm{trough}}=\frac{\log\Gamma}{\rho},
\end{equation}
the first random and slowly moving, the second a fixed constant. Section
One's soft-wall reduction applies with its fixed wall replaced by
$\max$ of these.
\end{lemma}

\begin{proof}
Solve $B_te^{\,b-x}=1$ and $\Gamma e^{-\rho x}=1$.
\end{proof}

\begin{theorem}[Exact absorption of a linear background; \textbf{P}]
\label{tc:thm:absorb}
Suppose $\log B_u=a+cu$ exactly --- the Gompertz phase, $a=\log n$,
$c=\epsilon^{2}/2$. Then persistence above the crest wall
\eqref{tc:eq:killwalls} is \emph{identical} to Section~\ref{sh:sec}'s fixed-wall
problem with initial data
\begin{equation}\label{tc:eq:shift}
\mu_0\longmapsto\mu_0-c,\qquad D_0\longmapsto D_0-a-b .
\end{equation}
By that section's gauge proposition the background therefore moves
prefactors and timescales alone, and no asymptotic exponent.
\end{theorem}

\begin{proof}
With $D_u=D_0+\mu_0u+\sigma I_u$, $I_u=\int_0^u W_s\,ds$, the constraint
$D_u\ge a+b+cu-x_0$ reads
$\sigma\bigl(I_u-\tfrac{c-\mu_0}{\sigma}u\bigr)\ge a+b-x_0-D_0$, and
$I_u-\kappa u=\int_0^u(W_s-\kappa)ds$ is the integral of a Brownian
motion started at $-\kappa$: the same Kolmogorov diffusion from shifted
data.
\end{proof}

\begin{remark}[A headwind]\label{tc:rem:headwind}
An exponentially growing background is a \emph{constant headwind} of
size $\epsilon^{2}/2$ on the improvement rate: the system must improve
at that rate merely to hold its hazard fixed, and its initial credit is
debited by $\log n+b$. Ageing, here, is not a new mechanism but an
offset.
\end{remark}

\begin{theorem}[Sublinear walls are invisible; \textbf{A}]
\label{tc:thm:sublinear}
If both walls are $o(u^{3/2})$ almost surely --- which holds here, the
crest wall being $O(u)$ in the delocalised phase and $O(\sqrt u)$ in the
frozen phase, and the trough wall constant --- then, granting the
moving-boundary robustness of Section~\ref{sh:sec}'s Ansatz class, the
persistence exponent is $\tfrac14$, unchanged.
\end{theorem}

\begin{proof}[Sketch and status]
After rescaling by $t^{3/2}$ the boundary converges to zero uniformly on
compacts, so the constraint converges to the fixed-wall constraint.
Converting this to an exponent identity is the moving-boundary
robustness discussed in Section~\ref{sh:sec} and surveyed in \cite{AS15}.
\end{proof}

% ---------------------------------------------------------------------
\subsection{Strength and anatomy as coordinates}\label{tc:sub:polar}

The author's requirement is that strength and anatomy be a two-part
parametrisation which, when both are random, is \emph{equivalent} to the
one-parameter parametrisation of $x$ by $\mu$ --- that is, a change of
coordinates on the state, introducing nothing new. This exists, is
canonical, and is available adaptedly.

\subsubsection{The decomposition}\label{tc:subsub:polar}

\begin{definition}[Polar coordinates of the state]\label{tc:def:polar}
Assume $x_0=\mu_0=0$. Normalise by Section~\ref{sh:sec}'s own scaling exponents
and orthogonalise the velocity against the position:
\begin{equation}\label{tc:eq:polarcoords}
\hat x_t=\frac{x_t}{\sigma t^{3/2}},
\qquad
Q_t=\frac{2\mu_t t-3x_t}{\sqrt3\,\sigma t^{3/2}},
\end{equation}
and set $R_t=\sqrt{\hat x_t^{2}+Q_t^{2}}$,
$\Theta_t=\arg(\hat x_t+iQ_t)$, with \emph{amplitude}
$\Ac_t:=\sigma t^{3/2}R_t$. Then
\begin{equation}\label{tc:eq:polar}
\boxed{\;x_t=\sigma\,t^{3/2}R_t\cos\Theta_t=\Ac_t\cos\Theta_t\;}
\end{equation}
We name $R$ (equivalently $\log R$) the \emph{strength} and $\Theta$ the
\emph{anatomy}.
\end{definition}

\begin{theorem}[Properties of the decomposition; \textbf{P}]
\label{tc:thm:polar}
Under the hypotheses of Definition~\ref{tc:def:polar}:
\begin{enumerate}[label=\emph{(\roman*)},itemsep=1pt]
\item \emph{Exact and bijective.} \eqref{tc:eq:polar} is an identity,
and $(\mu_t,x_t)\leftrightarrow(R_t,\Theta_t)$ is a bijection at each
$t>0$: the pair carries exactly the information of the state, no more
and no less.
\item \emph{Adapted.} $(R_t,\Theta_t)$ is a function of
$(\mu_t,x_t,t)$ alone --- no Hilbert transform and no lookahead. This is
why the velocity is orthogonalised against the position rather than the
textbook Rice envelope \cite{Ric44,Ric45,CL67} being used.
\item \emph{Isotropic Gaussian.} $\Cov(\hat x_t,Q_t)=0$ and
$\Var\hat x_t=\Var Q_t=\tfrac13$ for every $t$.
\item \emph{Marginal laws.} $R_t\sim\mathrm{Rayleigh}(1/\sqrt3)$ and
$\Theta_t\sim\mathrm{Uniform}[0,2\pi)$, independent at each fixed $t$.
\item \emph{Positivity and sign.} $R_t>0$, so $R=e^{\alpha}$ with
$\alpha=\log R$ automatically; and $x_t<0\iff\cos\Theta_t<0$.
\item \emph{Log-time stationarity.} $(\hat x_{e^{s}},Q_{e^{s}})_{s}$ is
a stationary Gaussian process, hence $(R,\Theta)$ are jointly stationary
in log-time and ergodic.
\end{enumerate}
\end{theorem}

\begin{proof}
(i) is algebra, the map being $(\mu,x)\mapsto(\hat x,Q)$ linear and
invertible for $t>0$ followed by planar polar coordinates. (ii) is
inspection of \eqref{tc:eq:polarcoords}.

(iii) With $\sigma=1$, $x_t=I_t$ and $\mu_t=W_t$, so
$\Var I_t=t^{3}/3$, $\Cov(W_t,I_t)=t^{2}/2$, $\Var W_t=t$. The
regression coefficient of $W_t$ on $I_t$ is
$(t^{2}/2)/(t^{3}/3)=3/2t$, so $W_t-\tfrac{3}{2t}I_t$ is orthogonal to
$I_t$; multiplying by $2t$ gives $2tW_t-3I_t$, with variance
$4t^{3}-6t^{3}+3t^{3}=t^{3}$. Dividing by $\sqrt3\,t^{3/2}$ gives
variance $\tfrac13$, matching $\Var\hat x_t=\tfrac13$. The $\sqrt3$ is
therefore not a tuning constant but exactly what equalises the two
variances.

(iv) An isotropic centred Gaussian pair with common variance $s^{2}$ has
Rayleigh modulus with parameter $s$ and uniform independent argument;
here $s=1/\sqrt3$. (v) is immediate. (vi) Both components are
$3/2$-self-similar functionals, so each is a Lamperti transform;
stationarity of the pair follows from stationarity of the joint
covariance under $t\mapsto ct$.
\end{proof}

\subsubsection{What independence does and does not mean}
\label{tc:subsub:indep}

The marginal independence in Theorem~\ref{tc:thm:polar}(iv) is a
statement at each fixed time. Whether it survives as a statement about
the \emph{processes} is a separate question, and it is decidable.

\begin{proposition}[The polar processes are dependent; \textbf{P}]
\label{tc:prop:notindep}
Normalise $\sigma=1$ and take $1\le t$. Then
\begin{align}
\Cov(\hat x_1,\hat x_t)&=\frac{3t-1}{6t^{3/2}},
&\Cov(Q_1,Q_t)&=\frac{3-t}{6t^{3/2}},
\label{tc:eq:covs}\\
\Cov(\hat x_1,Q_t)&=\frac{\sqrt3\,(1-t)}{6t^{3/2}},
&\Cov(Q_1,\hat x_t)&=\frac{\sqrt3\,(t-1)}{6t^{3/2}} .
\notag
\end{align}
Consequently:
\begin{enumerate}[label=\emph{(\roman*)},itemsep=1pt]
\item \emph{Exact antisymmetry.}
$\Cov(\hat x_1,Q_t)+\Cov(Q_1,\hat x_t)=0$ identically --- one half of
the rotation-invariance condition holds exactly.
\item \emph{Failure of rotation invariance.} The other half fails:
\begin{equation}\label{tc:eq:deviation}
\Cov(\hat x_1,\hat x_t)-\Cov(Q_1,Q_t)=\frac{2(t-1)}{3t^{3/2}}
\;\not\equiv\;0 .
\end{equation}
\item Hence $(R_t)$ and $(\Theta_t)$ are \emph{not} independent as
processes.
\item \emph{Where the dependence lives.} In log-lag $s=\log t$ the
deviation \eqref{tc:eq:deviation} vanishes at $s=0$ and as
$s\to\infty$, and is maximised at $s=\log3$, where it equals
$4/(3\cdot3^{3/2})\approx0.2566$ against the common variance
$\tfrac13$. The dependence is therefore substantial at intermediate
log-lags and negligible at short and long ones.
\end{enumerate}
\end{proposition}

\begin{proof}
The covariances follow from $\E[I_1I_t]=t/2-\tfrac16$,
$\E[W_1I_t]=t-\tfrac12$, $\E[I_1W_t]=\tfrac12$, $\E[W_1W_t]=1$ (each by
Fubini on $\E[W_aW_b]=a\wedge b$) inserted into
\eqref{tc:eq:polarcoords}. (i) and (ii) are then immediate.
(iii): for a stationary planar Gaussian process, independence of modulus
and argument processes requires invariance of the law under joint
rotation, which requires both equality of the diagonal covariances and
antisymmetry of the off-diagonal ones; (ii) fails the first. (iv):
substituting $t=e^{s}$ in \eqref{tc:eq:deviation} and differentiating,
$\frac{d}{ds}\bigl[(e^{s}-1)e^{-3s/2}\bigr]
=e^{-3s/2}\bigl(\tfrac32-\tfrac12e^{s}\bigr)$, which vanishes at
$e^{s}=3$.
\end{proof}

\begin{remark}[The honest form of ``separate knobs'']
\label{tc:rem:knobs}
Strength and anatomy are genuinely separate knobs \emph{in the sense
that matters for describing the state at a given moment}: at each fixed
time they are independent, with explicit laws, and together they carry
the whole state. They are not separate in the sense of evolving
independently. A system whose strength is currently large is more likely
to have had a particular anatomy a factor of $3$ ago in time. The right
summary is that the decomposition factorises the \emph{marginal}
geometry exactly and the \emph{dynamics} not at all --- which is
unsurprising, since the dynamics is precisely what the single driver
$\mu$ supplies.
\end{remark}

\subsubsection{Anatomy as phase, and an exact half}
\label{tc:subsub:half}

\begin{remark}[Reading]\label{tc:rem:reading}
$R$ is how much muffling structure exists; $\Theta$ is which way it is
pointed. Bad anatomy is not less material but the same material arranged
wrongly. Concretely $\Theta\approx0$ is the crest of a muffling
excursion, $\Theta\approx\pi$ the trough, $\Theta\approx\pm\pi/2$ a zero
crossing: \emph{anatomy is phase within the excursion cycle}. Note that
this is a different sense of ``trough'' from the trough channel of
\eqref{tc:eq:hazard} --- the two coincide only when the amplitude is
large enough to carry $x$ past the wall, which is the content of
\S\ref{tc:sub:room}.
\end{remark}

\begin{theorem}[Bad anatomy occupies exactly half of logarithmic time;
\textbf{P}]\label{tc:thm:half}
Almost surely,
\begin{equation}\label{tc:eq:half}
\lim_{S\to\infty}\frac1S\,\mathrm{Leb}
\bigl\{s\in[0,S]:\cos\Theta_{e^{s}}<0\bigr\}=\frac12 .
\end{equation}
Equivalently, the set $\{t:x_t<0\}$ has logarithmic density exactly
$\tfrac12$.
\end{theorem}

\begin{proof}
$\cos\Theta_t<0\iff\hat x_t<0\iff x_t<0$. The process
$Z_s:=\hat x_{e^{s}}$ is stationary, centred Gaussian and ergodic
(Section~\ref{sh:sec}'s Lamperti argument), so Birkhoff gives density
$\PP(Z_0<0)=\tfrac12$ by symmetry.
\end{proof}

\begin{remark}[What the sharpening owes to the polar form]
\label{tc:rem:owes}
Theorem~\ref{tc:thm:half} genuinely sharpens Section~\ref{sh:sec}, whose
no-viability proof used only the qualitative input
$\PP(Z_0<-\varepsilon)>0$; here the constant is exact. But the
sharpening is owed to \emph{symmetry and ergodicity of $Z$}, not to the
polar decomposition: $\{\cos\Theta<0\}=\{x<0\}$, and the argument never
uses $R$. The polar form makes the statement transparent --- it exhibits
$\tfrac12$ as the uniform measure of a half-circle --- and it is right
to state it here, but it is not required. We flag this because the
converse impression is easy to form.
\end{remark}

% ---------------------------------------------------------------------
\subsection{The waiting room}\label{tc:sub:room}

\subsubsection{Two kinds of wall}\label{tc:subsub:twowalls}

A distinction must be drawn that the informal account elides. There are
two different senses in which a channel ``has a wall'', and both are
needed.

\begin{definition}[Shape and killing walls]\label{tc:def:walls}
The \emph{shape walls} are where a channel rises above the background
floor, i.e.\ where it becomes comparable to $B_t$:
\begin{equation}\label{tc:eq:shapewalls}
x^{\mathrm{crest}}_{\mathrm{shape}}\approx b,
\qquad
x^{\mathrm{trough}}_{\mathrm{shape}}\approx-\frac{k_t}{\rho},
\qquad
k_t:=\log\frac{\E[B_t]}{\Gamma},
\end{equation}
so that the room is the \emph{asymmetric} interval $[-k_t/\rho,\,b]$,
containing the origin in its interior.
The \emph{killing walls} are where a channel makes $\Lambda$ accumulate,
i.e.\ where it becomes $O(1)$ in absolute terms; these are
\eqref{tc:eq:killwalls}. The shape walls govern what the hazard
\emph{looks like} as a function of $x$; the killing walls govern
Section~\ref{sh:sec}'s persistence geometry. They are different objects and are
not interchangeable.
\end{definition}

\begin{remark}[$\Gamma$ primitive, $k$ derived]\label{tc:rem:kderived}
$\Gamma$ is a primitive of Definition~\ref{tc:def:model}; the width
$k_t$ is \emph{not} a parameter but a bookkeeping quantity recording
where $\Gamma$ sits relative to the background. Two consequences follow
immediately and are worth separating, since they run in opposite
directions:
\begin{itemize}[itemsep=1pt]
\item \emph{The room widens.} In the delocalised phase
$\E[B_t]=ne^{\epsilon^{2}t/2}$, so $k_t=\log(n/\Gamma)+\epsilon^{2}t/2$
grows linearly and the shape-wall separation $k_t/\rho$ widens at rate
$\epsilon^{2}/2\rho$: \emph{grace grows as the background rises}.
\item \emph{The killing walls separate faster.} The crest killing wall
$\log(2B_t)$ rises at $\epsilon^{2}/2$ while the trough killing wall
$\rho^{-1}\log\Gamma$ is constant, so their separation grows at
$\epsilon^{2}/2$.
\end{itemize}
Both are harmless asymptotically, being $O(t)=o(t^{3/2})$.
\end{remark}

\subsubsection{The void, and when it is flat}\label{tc:subsub:void}

\begin{proposition}[The room, and when it is flat; \textbf{P}]
\label{tc:prop:flat}
The crest is a sigmoid of transition width $O(1)$ centred at $x=b$:
$\varsigma_b\approx1$ for $x\lesssim b-2$ and $\varsigma_b\approx
e^{\,b-x}$ for $x\gtrsim b+2$. Hence the total hazard is flat, at
$\approx B_t$, on
\begin{equation}\label{tc:eq:flat}
-\frac{k_t}{\rho}\;\ll\;x\;\lesssim\;b-2 ,
\end{equation}
a genuine plateau when $k_t/\rho+b\gg1$. It contains $x=0$ in its
interior precisely when $b\gtrsim2$, which is the structural role of the
offset.
\end{proposition}

\begin{proof}
$\varsigma_b(x)\in(0.88,1.05)\cdot$ its plateau value for
$x\le b-2$, and the trough exceeds the crest below $-k_t/\rho$ by
Proposition~\ref{tc:prop:dropout}(i).
\end{proof}

\begin{constraint}[Internal onsets before external]\label{tc:C5}
As the amplitude grows, $x$ reaches the trough wall before the crest
wall if and only if
\begin{equation}\label{tc:eq:C5}
\boxed{\;b\;>\;\frac{k_t}{\rho}\;}
\end{equation}
--- the offset must exceed the trough's half-width. Under
\eqref{tc:eq:C5} a system of unfavourable anatomy suffers internal
damage first, and only later does its external hazard begin to fall.
\end{constraint}

\begin{remark}[Why the offset is needed]\label{tc:rem:whyoffset}
At $b=0$ the origin sits at the crest's inflection, where the local gain
is $\tfrac12$: a system with no muffling at all would already be half
effective, contradicting the reading of the void as the regime in which
the system contributes nothing. The offset moves the origin into the
plateau, where by Theorem~\ref{tc:thm:gain}(iii) the gain is $O(e^{-b})$.
The void then contains its own centre, and \eqref{tc:eq:C5} orders the
two onsets.
\end{remark}

\begin{remark}[Void as mutedness]\label{tc:rem:mutedness}
The void is where the system contributes nothing. That is stated
qualitatively by Remark~\ref{tc:rem:regimes} and quantitatively by
Proposition~\ref{tc:prop:mingain}: the gain $\Psi'$, which measures how
much the driver moves the hazard, falls to
$\kappa(\rho)e^{-(\rho b+k)/(1+\rho)}$ on the void, exponentially small
in the room's width. A system in the waiting room is not merely at low hazard;
its efforts have almost no purchase, and the wider the room the less.
\end{remark}

\subsubsection{Activation: time supplies the interpolation}
\label{tc:subsub:activation}

\begin{proposition}[Activation time; \textbf{S}]\label{tc:prop:tact}
In log-time $s=\log t$,
\begin{equation}\label{tc:eq:logamp}
\log\Ac_{e^{s}}=\log\sigma+\tfrac32s+\log R_{e^{s}},
\qquad R\ \text{stationary},
\end{equation}
so the amplitude is exponentially small at early log-times and grows
exponentially, with a stationary fluctuation riding on a deterministic
ramp of slope $\tfrac32$. The void therefore ends, and the channels
engage, at
\begin{equation}\label{tc:eq:tact}
t^{\mathrm{int}}_{\mathrm{act}}\asymp\Bigl(\frac{k}{\rho\,\sigma}
\Bigr)^{2/3},
\qquad
t^{\mathrm{ext}}_{\mathrm{act}}\asymp\Bigl(\frac{b}{\sigma}\Bigr)^{2/3},
\end{equation}
the times at which the typical amplitude reaches the trough and crest
walls respectively. Under Constraint~\ref{tc:C5} these are ordered,
$t^{\mathrm{int}}_{\mathrm{act}}<t^{\mathrm{ext}}_{\mathrm{act}}$: the
internal channel engages first and the external one only later.
\end{proposition}

\begin{proof}[Scaling argument]
$\Ac_t=\sigma t^{3/2}R_t$ with $R_t=O_P(1)$ stationary; setting
$\sigma t^{3/2}\asymp k/\rho$ and $\sigma t^{3/2}\asymp b$ give
\eqref{tc:eq:tact}. Dimensions:
$[(k/\rho\sigma)^{2/3}]=(\mathsf T^{3/2})^{2/3}=\mathsf T$. A sharp
constant would require the first-passage law of the amplitude, which we
do not compute.
\end{proof}

\begin{remark}[Why no capacity parameter is needed]
\label{tc:rem:nocapacity}
Revision~1 installed a capacity (strength) parameter with a leak, whose
purpose was to produce a default regime in which the system contributes
nothing and $\mu$ matters only rarely. Equation \eqref{tc:eq:logamp}
supplies that regime for free: the amplitude of $x$ is exponentially
small in log-time at early times, so the void is not an assumption to be
installed but \emph{what the model already does}, and it is left at
\eqref{tc:eq:tact} --- in two stages, internal then external --- without
any additional primitive. The interpolation
parameter is $k$ --- large $k$, long void, background-dominated; small
$k$, immediate engagement --- and $k$ is itself derived
(Remark~\ref{tc:rem:kderived}).
\end{remark}

\begin{remark}[The cost of dropping capacity, recorded]
\label{tc:rem:cost}
This is a real loss and should not be glossed.
\begin{center}\small
\begin{tabular}{@{}lll@{}}
\toprule
& capacity (revision~1, dropped) & $k$ (retained)\\
\midrule
what it moves & an \emph{exponent}, $\tfrac14\to\tfrac12$ &
a \emph{timescale}\\
sharpness & genuine threshold at $S_c$ & crossover at
$t_{\mathrm{act}}$\\
\bottomrule
\end{tabular}
\end{center}
Revision~1's sharp-threshold theorem is therefore gone, and its ``sharp
in parameter space'' register weakens to a crossover. What replaces it
is the threshold at $\rho=1$
(Proposition~\ref{tc:prop:dropout}(ii)), beyond which the background
still reaches the survivor's hazard and below which the system closes
off from the external world. That is a threshold in a structural
parameter with a mechanical meaning rather than one installed to produce
a threshold --- but, as Remark~\ref{tc:rem:rhorole}(c) records, it moves
no exponent. The exchange is favourable and honest, not free.
\end{remark}

% ---------------------------------------------------------------------
\subsection{The kinked gain, and what the driver was concealing}
\label{tc:sub:gain}

\subsubsection{The exact gain}\label{tc:subsub:gain}

\begin{theorem}[Soft-kink gain; \textbf{P}]\label{tc:thm:gain}
Fix $B_t\equiv B$ constant and define the \emph{gain}
\begin{equation}\label{tc:eq:Psi}
\Psi(x):=-\log\frac{h^{\mathrm{tot}}(x)}{h^{\mathrm{tot}}(0)},
\qquad\text{so that}\qquad
h^{\mathrm{tot}}(x)=h^{\mathrm{tot}}(0)\,e^{-\Psi(x)}
\end{equation}
identically, with $h^{\mathrm{tot}}(0)=B+\Gamma$. Then:
\begin{enumerate}[label=\emph{(\roman*)},itemsep=1pt]
\item $\Psi(0)=0$ and $\Psi$ is smooth and strictly increasing;
\item $\Psi(x)=x+O(e^{-x})$ as $x\to+\infty$ and
$\Psi(x)=\rho x+O(1)$ as $x\to-\infty$: a \emph{soft kink} of slopes
$1$ and $\rho$;
\item the local gain is a hazard-weighted average of the two channels'
slopes,
\begin{equation}\label{tc:eq:Psiprime}
\Psi'(x)=\frac{h^{\mathrm{crest}}(x)\,\frac{e^{x}}{1+e^{x}}
+h^{\mathrm{trough}}(x)\,\rho}
{h^{\mathrm{crest}}(x)+h^{\mathrm{trough}}(x)}\;\in\;(0,\rho),
\end{equation}
with $\Psi'(0)\approx e^{-b}+\rho e^{-k}$, $\Psi'\to1$ as
$x\to+\infty$, $\Psi'\to\rho$ as $x\to-\infty$, and an interior minimum
on the void quantified in Proposition~\ref{tc:prop:mingain}.
\end{enumerate}
Consequently, with constant background, the model is exactly Section
One's with the credit passed through $\Psi$: Constraint~\ref{tc:C4} is
met.
\end{theorem}

\begin{proof}
(i) is immediate from \eqref{tc:eq:Psi} and positivity of
$h^{\mathrm{tot}}$; monotonicity from (iii). (ii): as $x\to+\infty$,
$h^{\mathrm{tot}}=2Be^{-x}(1+O(e^{-(\rho-1)x}))$, so $\Psi=x+O(\cdot)$;
as $x\to-\infty$, $h^{\mathrm{tot}}=\Gamma e^{-\rho x}(1+O(e^{(\rho-1)
x}))$, so $\Psi=\rho x+O(1)$. (iii): for a sum $h=h_1+h_2$ one has
$-h'/h=(h_1s_1+h_2s_2)/(h_1+h_2)$ with $s_i=-h_i'/h_i$; here
$s_2=\rho$ and, differentiating $2B/(1+e^{x})$,
$s_1=e^{x-b}/(1+e^{x-b})\in(0,1)$. The stated limits follow by
evaluating the weights, using $\Gamma/B=e^{-k}$ at $x=0$.
\end{proof}

\begin{proposition}[The void is muted, exponentially in $k$;
\textbf{P}]\label{tc:prop:mingain}
On the void the gain \eqref{tc:eq:Psiprime} attains an interior minimum
at
\begin{equation}\label{tc:eq:xstar}
x^{*}=\frac{b-k+2\log\rho}{1+\rho},
\qquad
\min\Psi'=\kappa(\rho)\,e^{-\frac{\rho b+k}{1+\rho}},
\qquad
\kappa(\rho):=\rho^{\frac{2}{1+\rho}}+\rho^{\frac{1-\rho}{1+\rho}} .
\end{equation}
The gain therefore does not merely dip but is \emph{exponentially small
in $(\rho b+k)/(1+\rho)$}: the wider the room --- on either side --- the
more completely the driver is muted within it. For $\rho=\tfrac32$,
$\kappa=2.045$; for $\rho=\tfrac52$, $\kappa=2.443$.
\end{proposition}

\begin{proof}
On the void the crest is saturated, so its weight is $\approx1$ and its
slope is $s_{1}\approx e^{x-b}$, while the trough weight is
$\approx(\Gamma/B)e^{-\rho x}=e^{-k}e^{-\rho x}$. Hence
$\Psi'\approx f(x):=e^{x-b}+\rho e^{-k}e^{-\rho x}$, and $f'(x)=0$ gives
$e^{(1+\rho)x}=\rho^{2}e^{\,b-k}$, which is \eqref{tc:eq:xstar}.
Substituting back, both terms carry the common factor
$e^{-(\rho b+k)/(1+\rho)}$ with coefficients $\rho^{2/(1+\rho)}$ and
$\rho^{(1-\rho)/(1+\rho)}$, summing to the stated value. (The
approximation is verified numerically to four significant figures over
$\rho\in\{1.5,2.5\}$, $b\in[6,14]$, $k\in[8,20]$.)
\end{proof}

\begin{remark}[Correction to the reference specification]
\label{tc:rem:speccorr}
Two points where the informal specification requires amendment, both
recorded in the ledger.

\emph{(a)} The closed form $\psi_\rho(x)=-\log(e^{-x}+e^{-\rho x})$ is
\emph{not} an identity for \eqref{tc:eq:hazard}: at $x=0$ the model
gives $B+\Gamma$ while $\psi_\rho$ gives $2h_0$. It is the clean
representative of the soft-kink class, not the model's own gain;
\eqref{tc:eq:Psi} is exact by construction and has the same asymptotic
slopes, so nothing structural is lost.

\emph{(b)} The gain range is $(0,\rho)$ and not $[1,\rho]$. The crest
channel's local slope is $e^{x}/(1+e^{x})$, which tends to $0$, not $1$,
as $x\to-\infty$. This is not a blemish but the sharpest available
statement of what the void \emph{is}: on the plateau the gain falls to
$\kappa(\rho)e^{-(\rho b+k)/(1+\rho)}$
(Proposition~\ref{tc:prop:mingain}), so
the driver is muted and its efforts scarcely move the hazard. The gain
profile thus encodes all three regimes --- exponentially small in the
void, $\to1$ under shielding, $\to\rho$ under damage. The mutedness is a
property of a \emph{wide} room: if $\rho b+k$ is $O(1)$ the minimum is
$O(1)$ and there is no muted regime, consistently with
Proposition~\ref{tc:prop:flat}.
\end{remark}

\subsubsection{What the original driver was concealing}
\label{tc:subsub:concealed}

\begin{corollary}[State-dependent gain; \textbf{P}]
\label{tc:cor:concealed}
Along a path, the effective negative logarithmic derivative of the total
hazard is
\begin{equation}\label{tc:eq:effmu}
\tilde\mu_t=-\frac{d}{dt}\log h^{\mathrm{tot}}_t
=\Psi'(x_t)\,\mu_t,
\qquad \Psi'\in(0,\rho).
\end{equation}
That is: \emph{Section~\ref{sh:sec}'s driver was the present driver times a
state-dependent gain}, and the old Brownian motion was absorbing that
gain as extra noise. The revision therefore extracts structure from the
original noise term rather than adding a primitive to it.
\end{corollary}

\begin{proof}
Chain rule on \eqref{tc:eq:Psi} with $\dot x=\mu$.
\end{proof}

\begin{remark}[Reading of the gain]\label{tc:rem:gainreading}
The new $\mu$ is correspondingly narrower than the old, and the
difference is now explicit rather than hidden. Where the old model
attributed all variation in the log-hazard's slope to the driver, the
present one attributes part of it to \emph{where the system is}: the
same effort buys a full unit of log-hazard reduction under shielding,
$\rho$ units of damage in a trough, and nothing at all in the void.
\end{remark}

% ---------------------------------------------------------------------
\subsection{Asymptotics: Section~\ref{sh:sec} survives intact}
\label{tc:sub:asymptotics}

\begin{theorem}[Section~\ref{sh:sec} is undisturbed; \textbf{A}]
\label{tc:thm:survives}
In the two-channel model, for every $\rho>0$, $\Gamma>0$ and $b\ge0$:
\begin{enumerate}[label=\emph{(\roman*)},itemsep=1pt]
\item the persistence exponent is $\tfrac14$, exactly;
\item $\PP(T>t)\asymp Ct^{-1/4}$;
\item $\E[T]=\infty$ while $T<\infty$ a.s.\ (null recurrence);
\item the survivorship interpolation of Section~\ref{sh:sec}, its pure Doob
limit, its $s/4t$ linear response and its critical mode, are unchanged.
\end{enumerate}
\emph{No truncation window and no $\gamma=0$ caveat}:
Constraint~\ref{tc:C3} holds as a theorem.
\end{theorem}

\begin{proof}
By Lemma~\ref{tc:lem:walls} the killing walls are $\log(2B_t)$ and
$\rho^{-1}\log\Gamma$. The latter is constant, hence an $O(1)$ shift,
which Section~\ref{sh:sec} proved to be gauge. The former is $O(t)$ in the
delocalised phase and $O(\sqrt t)$ in the frozen phase, hence
$o(t^{3/2})$, and is absorbed exactly in the linear case by
Theorem~\ref{tc:thm:absorb} and invisibly in general by
Theorem~\ref{tc:thm:sublinear}; the offset $b$ enters
\eqref{tc:eq:shift} as a further $O(1)$ shift and is likewise gauge.
Survivors live at large positive $x$, where by
Proposition~\ref{tc:prop:dropout}(ii) the dominant channel is the crest
if $\rho>1$ and the trough if $\rho<1$; in either case the dominant
channel is an exponential in $x$, so the killing wall is a level of the
credit and the persistence problem is Section~\ref{sh:sec}'s up to a constant
rescaling. The conclusions are then Section~\ref{sh:sec}'s, transported --- and,
notably, for every $\rho>0$ rather than only $\rho>1$.
\end{proof}

\begin{remark}[This repairs revision~1's defect]\label{tc:rem:repair}
Revision~1's additive internal hazard contributed $\gamma t$ to
$\Lambda$, destroying the $t^{-1/4}$ tail and forcing a truncation
window with finite life expectancy $\asymp g^{-3/4}$. The trough channel
does not, and the reason is structural rather than fortunate:
$\Gamma e^{-\rho x}$ is small exactly where survivors live, whereas a
constant is not small anywhere. Both walls are $O(1)$-or-gauge shifts,
and Section~\ref{sh:sec} proved such shifts to be gauge.
\end{remark}

\begin{remark}[What was given up]\label{tc:rem:givenup}
Both channels are decreasing in $x$, so the total hazard is monotone and
\emph{more strength is always better}. Revision~1's U-shaped hazard, and
with it the optimal finite strength and the $(S,A)$ optimisation, are
gone. This is the price of Constraint~\ref{tc:C3}: an interior optimum
requires a term that grows with $x$, and any such term destroys the
tail. An optimum and an undisturbed Section~\ref{sh:sec} were never going to
coexist. What survives of the earlier optimisation story is the
threshold at $\rho=1$ and the trade recorded in
Remark~\ref{tc:rem:cost}.
\end{remark}

% ---------------------------------------------------------------------
\subsection{Transfer: how much of Section~\ref{sh:sec} survives, and where the
differences go}\label{tc:sub:transfer}

Theorem~\ref{tc:thm:survives} asserted that Section~\ref{sh:sec}'s conclusions
carry over. This subsection makes the transfer quantitative: it
enumerates the four ways the two-channel hazard differs from
$h_{0}e^{-x}$, shows that each lands in gauge or in the prefactor, and
isolates the one place where the two-channel model is genuinely
\emph{better} conditioned than Section~\ref{sh:sec}.

\subsubsection{The four differences}\label{tc:subsub:four}

\begin{center}\small
\begin{tabular}{@{}llll@{}}
\toprule
difference & size & where it goes & reference\\
\midrule
wall location $\log B_u+b$ & $\Theta(u)$ & gauge (exact absorption) &
Thm.~\ref{tc:thm:absorb}\\
graded boundary of width $W_u$ & $\Theta(u)$ & gauge (collapses under
rescaling) & Thm.~\ref{tc:thm:collapse}\\
waiting-room passage & $O(1)$ & prefactor & Prop.~\ref{tc:prop:roomcost}\\
background fluctuation & $O(\sqrt u)$ & gauge &
Thm.~\ref{tc:thm:sublinear}\\
\bottomrule
\end{tabular}
\end{center}

\noindent Nothing reaches an exponent. The transfer is therefore exact
at the level at which Section~\ref{sh:sec} states its results.

\begin{lemma}[Above the room the model is Section~\ref{sh:sec}'s; \textbf{P}]
\label{tc:lem:aboveroom}
For $x\ge b$,
\begin{equation}\label{tc:eq:sandwich}
\tfrac12\bigl(1+e^{-b}\bigr)e^{\,b-x}\;\le\;\varsigma_b(x)\;\le\;
\bigl(1+e^{-b}\bigr)e^{\,b-x},
\end{equation}
so $h^{\mathrm{crest}}\asymp B_ue^{\,b-x}$ with constants in
$[\tfrac12,1]\cdot(1+e^{-b})$. If moreover $\rho>1$ then
$h^{\mathrm{trough}}/h^{\mathrm{crest}}\to0$ as $x\to\infty$, so
$h^{\mathrm{tot}}\asymp B_ue^{\,b-x}$: above the room the hazard is
Section~\ref{sh:sec}'s $h_{0}e^{-x}$ with $h_{0}=B_ue^{b}$, up to bounded
multiplicative constants.
\end{lemma}

\begin{proof}
For $x\ge b$ one has $1\le e^{x-b}$, hence
$e^{x-b}\le1+e^{x-b}\le2e^{x-b}$; invert. The trough comparison is
Proposition~\ref{tc:prop:dropout}(ii).
\end{proof}

\subsubsection{No viability, by a different route}\label{tc:subsub:noviab}

Section~\ref{sh:sec}'s no-viability theorem was proved by exhibiting an integrand
that \emph{blows up}: on a set of infinite measure the credit satisfies
$D_u\le-\varepsilon u^{3/2}$, whence $e^{-D_u}\ge
e^{\sigma\varepsilon u^{3/2}}\to\infty$. That argument is unavailable
here whenever the crest channel acts alone, since $\varsigma_b$ is
\emph{bounded}. The conclusion nevertheless survives, by a weaker and
more robust mechanism: a positive floor on an infinite set.

\begin{theorem}[No viability in the two-channel model; \textbf{P}]
\label{tc:thm:noviab2}
For every $\sigma,\epsilon>0$, $n\ge1$, $b\ge0$, $\rho>0$ and every
$\Gamma\ge0$ --- including $\Gamma=0$, where the total hazard is
bounded --- one has $\Lambda_{\infty}=\infty$ almost surely. Hence
$\bar F_{\infty}=0$ a.s.\ and $T<\infty$ a.s.
\end{theorem}

\begin{proof}
Discard the trough, which is nonnegative. On $\{x_u\le0\}$ we have
$\varsigma_b(x_u)\ge\varsigma_b(0)=1$, so
\[
\Lambda_{t}\;\ge\;\int_{0}^{t}B_u\varsigma_b(x_u)\,du
\;\ge\;\Bigl(\inf_{u\le t}B_u\Bigr)\,
\mathrm{Leb}\{u\le t:x_u\le0\}.
\]
By Theorem~\ref{tc:thm:half} the set $\{u:x_u\le0\}$ has logarithmic
density $\tfrac12$; writing $u=e^{s}$,
$\mathrm{Leb}\{u\le t:x_u\le0\}=\int_{0}^{\log t}\one\{Z_s\le0\}e^{s}ds
\to\infty$ a.s., since the integrand is positive on a set of density
$\tfrac12$ and $e^{s}\to\infty$. As $\inf_{u\le t}B_u>0$ a.s.\ for each
$t$ and $B$ does not tend to zero, $\Lambda_{\infty}=\infty$.
\end{proof}

\begin{remark}[What the exact half is for]\label{tc:rem:halfuse}
Theorem~\ref{tc:thm:half} was presented in \S\ref{tc:sub:polar} as a
sharpening of Section~\ref{sh:sec}. It is more than that: it is the engine of the
proof just given. Where Section~\ref{sh:sec} needed an unbounded integrand,
the two-channel model needs only that the system spends a positive
fraction of logarithmic time at unfavourable anatomy --- and the polar
decomposition says that fraction is exactly one half, for every
parameter value. Boundedness of the hazard does not buy immortality; it
merely changes which argument rules it out.
\end{remark}

\subsubsection{The graded boundary collapses}\label{tc:subsub:collapse}

The substantive structural difference is that the killing boundary is
neither hard nor soft but \emph{graded}. Above $\log B_u+b$ the hazard
is negligible; across a band below it the hazard is $\asymp B_u$,
bounded; below $(\log\Gamma)/\rho$ it diverges.

\begin{theorem}[The band is asymptotically invisible; \textbf{A}]
\label{tc:thm:collapse}
Let
\begin{equation}\label{tc:eq:band}
W_u:=\underbrace{\log B_u+b}_{\text{crest wall}}
-\underbrace{\frac{\log\Gamma}{\rho}}_{\text{trough wall}}
=\Bigl(1-\frac1\rho\Bigr)\log B_u+b+\frac{k}{\rho}
\end{equation}
be the width of the graded band. In the delocalised phase
$W_u=\Theta(u)$ and in the frozen phase $W_u=\Theta(\sqrt u)$; in both,
\[
\frac{W_u}{u^{3/2}}\longrightarrow0 .
\]
Hence under Section~\ref{sh:sec}'s rescaling the whole band contracts to a single
point, and the graded boundary is asymptotically the hard wall of that
section. Granting the moving-boundary robustness of its Ansatz class,
the persistence exponent is $\tfrac14$.
\end{theorem}

\begin{proof}
Insert $\log B_u=\log n+\epsilon^{2}u/2$ (delocalised,
Corollary~\ref{tc:cor:gompertz}) or
$\log B_u=\epsilon\sqrt{2u\log n}$ (frozen,
Theorem~\ref{tc:thm:freezing}) into \eqref{tc:eq:band}; both are
$o(u^{3/2})$. The rescaling and the robustness step are as in
Theorem~\ref{tc:thm:sublinear}.
\end{proof}

\begin{remark}[Softening and re-hardening do not cancel each other ---
they both cancel against self-similarity]\label{tc:rem:netcancel}
It is tempting to read the situation as a competition: the logistic
bounds the hazard from above and would soften the wall, while the trough
is unbounded and re-hardens it, the two effects offsetting. That is not
what happens. The two effects act at \emph{different places} --- the
logistic at the top of the band, the trough at its bottom --- and
neither undoes the other. What renders both irrelevant is the third
thing: the band they enclose has width $o(u^{3/2})$, while the credit
they constrain fluctuates on scale $u^{3/2}$. A boundary structure of
any shape whatsoever, confined to a band of width $o(u^{3/2})$, is
invisible to the exponent. The net effect of the logistic and the trough
is therefore exactly nil asymptotically --- not by mutual cancellation
but because self-similarity outruns them both.
\end{remark}

\begin{proposition}[The tolerance shrinks: this model is better
conditioned than Section~\ref{sh:sec}'s; \textbf{P}]\label{tc:prop:tolerance}
Survival to time $t$ requires
\begin{equation}\label{tc:eq:tolerance}
\mathrm{Leb}\{u\le t:x_u\lesssim b\}\;=\;O\!\bigl(B_t^{-1}\bigr),
\end{equation}
since the hazard on the band is $\asymp B_u$ and $\Lambda_t=O(1)$ is
required. In the delocalised phase $B_t^{-1}=n^{-1}e^{-\epsilon^{2}t/2}
\to0$, so the admissible occupation below the wall tends to zero: the
problem converges to \emph{strict} persistence.
\end{proposition}

\begin{proof}
$\Lambda_t\ge\bigl(\inf_{u\le t}B_u\bigr)\cdot c\cdot
\mathrm{Leb}\{u\le t:x_u\lesssim b\}$ with $c=\varsigma_b(b)>0$;
survival to $t$ with probability bounded below forces $\Lambda_t=O(1)$.
\end{proof}

\begin{remark}[The Ansatz gap narrows]\label{tc:rem:gapnarrows}
Section~\ref{sh:sec}'s sole load-bearing heuristic was that persistence with
\emph{bounded occupation tolerance} shares the exponent of strict
persistence. The two-channel model does not evade that assumption, but
it weakens what is being assumed: by
Proposition~\ref{tc:prop:tolerance} the tolerance here is not $O(1)$ but
$O(B_t^{-1})\to0$. The richer, more ontologically committed model is
thus \emph{closer} to the rigorous case than the simpler one --- an
unusual direction for added structure to move a proof, and worth
recording as such.
\end{remark}

\subsubsection{The waiting room is a prefactor, and it is not a
shelter}\label{tc:subsub:roomcost}

\begin{proposition}[Cost of the room passage; \textbf{S}]
\label{tc:prop:roomcost}
The room $[-k/\rho,b]$ lies entirely below the crest killing wall
$\log B_u+b$, since $\log B_u>0$. A trajectory therefore traverses it
only before escaping upward, at time
$t^{\mathrm{ext}}_{\mathrm{act}}\asymp(b/\sigma)^{2/3}$
(Proposition~\ref{tc:prop:tact}), accumulating
\begin{equation}\label{tc:eq:roomcost}
\Lambda_{\mathrm{room}}\;\asymp\;
B_0\,t^{\mathrm{ext}}_{\mathrm{act}}
\;\asymp\;n\Bigl(\frac{b}{\sigma}\Bigr)^{2/3},
\end{equation}
an $O(1)$ quantity in $t$. Its entire effect is to multiply the
prefactor in $\PP(T>t)\asymp Ct^{-1/4}$ by
$e^{-\Theta\left(n(b/\sigma)^{2/3}\right)}$.
\end{proposition}

\begin{remark}[Grace, but not shelter]\label{tc:rem:notshelter}
The room deserves a more careful reading than ``grace interval''
suggests. Inside it the hazard is \emph{exactly the background}: neither
raised by damage nor lowered by shielding. It is therefore grace
relative to the trough --- unfavourable anatomy costs nothing while the
amplitude is small enough to keep $x$ above $-k/\rho$ --- but it is a
\emph{cost} relative to the crest, because survivors must be above the
wall and time spent in the room is time not spent shielding. Formula
\eqref{tc:eq:roomcost} prices that cost. The wider the room and the
larger the background, the more expensive the wait: the offset $b$ buys
the ordering of the two onsets (Constraint~\ref{tc:C5}) and pays for it
in the prefactor, at rate $n(b/\sigma)^{2/3}$ in the exponent.
\end{remark}

\subsubsection{The transfer, stated}\label{tc:subsub:transferstated}

\begin{theorem}[Transfer of Section~\ref{sh:sec}; \textbf{A}]
\label{tc:thm:transfer}
In the two-channel model, for every admissible parameter set,
\begin{equation}\label{tc:eq:transfer}
\PP(T>t)\;\asymp\;\underbrace{\hm\bigl(\mu_{0}-\tfrac{\epsilon^{2}}{2},
\;D_{0}-\log n-b\bigr)}_{\text{shifted initial data}}\;\cdot\;
\underbrace{e^{-\Theta\left(n(b/\sigma)^{2/3}\right)}}_{\text{room
passage}}\;\cdot\;t^{-1/4},
\end{equation}
with $\hm$ the persistence-harmonic function of Section~\ref{sh:sec}. In
consequence: $\E[T]=\infty$ while $T<\infty$ a.s.; residual life is
asymptotically Pareto$(\tfrac14)$; quantiles grow as the fourth power of
rarity; the effective population hazard is $1/4t$; and the survivorship
interpolation, its pure Doob limit, its $s/4t$ linear response and its
critical mode all hold verbatim, the limit being the same $Q$-process
computed at the shifted data.
\end{theorem}

\begin{proof}
Above the room the hazard is Section~\ref{sh:sec}'s up to bounded constants
(Lemma~\ref{tc:lem:aboveroom}); the wall is linear and absorbed exactly
into the shifted data (Theorem~\ref{tc:thm:absorb}, with the offset $b$
contributing a further constant); the band collapses
(Theorem~\ref{tc:thm:collapse}); the room contributes a bounded additive
term to $\Lambda$, hence a multiplicative prefactor
(Proposition~\ref{tc:prop:roomcost}). The listed consequences are those
of Section~\ref{sh:sec}, transported through
\eqref{tc:eq:transfer}.
\end{proof}

\begin{remark}[How close, in one sentence]\label{tc:rem:howclose}
Exactly as close as it is possible to be: every difference between the
two models is a shift of initial data or a bounded multiplicative
constant, and Section~\ref{sh:sec} proved both to be gauge. The two-channel model
buys its ontology --- a positive fluctuating background, a bounded
external channel obeying the causal asymmetry, an unbounded internal
channel, an ordered pair of onsets and a waiting room --- at a cost of
precisely zero asymptotic exponents, and it repays part of the price by
shrinking the one gap the earlier section had to assume.
\end{remark}

% ---------------------------------------------------------------------
\subsection{Numerical verification}\label{tc:sub:numerics}

Four tests specific to this section, additional to Section~\ref{sh:sec}'s.

\emph{(i) Freezing.} Simulate $B_t$ for $n\in\{10^{3},10^{4},10^{5}\}$
and plot $\log B_t/\log n$ against $\theta$; expect collapse onto
\eqref{tc:eq:freezing} with tangential contact at $\theta=\sqrt2$, and
sample-to-sample fluctuations $o(1)$ before the transition and $O(1)$
after, the frozen phase being non-self-averaging.

\emph{(ii) Polar marginals and their process dependence.} Verify that
$R_t$ is Rayleigh$(1/\sqrt3)$ and $\Theta_t$ uniform and independent at
fixed $t$ (a direct histogram check), and then verify
Proposition~\ref{tc:prop:notindep} by estimating
$\Cov(\hat x_1,\hat x_t)-\Cov(Q_1,Q_t)$ across log-lags: the prediction
is a positive bump vanishing at both ends with maximum at $t=3$, of
height $\approx0.257$. A flat zero would indicate an error in the
orthogonalisation.

\emph{(iii) Exponent robustness.} Estimate the survival tail for a grid
of $(\epsilon,n,\Gamma,\rho)$ and fit the local slope. Prediction:
$-\tfrac14$ throughout, with no dependence on any of the four ---
this is Theorem~\ref{tc:thm:survives} and the section's central
falsifiable claim.

\emph{(iv) The gain profile.} Compute $\Psi'$ from
\eqref{tc:eq:Psiprime} and check the limiting values of
Theorem~\ref{tc:thm:gain}(iii); then test
Proposition~\ref{tc:prop:mingain} by regressing $\log\min\Psi'$ on
$\rho b+k$ across a grid of $\rho$. The predicted slope is
$-1/(1+\rho)$, the
predicted intercept $\log\kappa(\rho)$, and the minimiser is at
\eqref{tc:eq:xstar}; all three are sharp, making this the most stringent
test in the section.

\emph{Falsifiers.} A fitted exponent differing from $-\tfrac14$, or
depending on $\Gamma$ or $\rho$ (contradicting
Theorem~\ref{tc:thm:survives}); a gain reaching a floor of $1$ rather
than $0$ on the plateau (contradicting
Remark~\ref{tc:rem:speccorr}(b)); process-level independence of $R$ and
$\Theta$ (contradicting Proposition~\ref{tc:prop:notindep}).

% ---------------------------------------------------------------------
\subsection*{Section summary}

The hazard has two channels. A bounded crest channel
$B_t\varsigma_b(x)$, a logistic shifted by an offset $b$, carries the
background, shields without limit but amplifies by at most $1+e^{-b}$,
and thereby enforces the causal asymmetry structurally rather than by
assumption. A trough channel $\Gamma e^{-\rho x}$ with brittleness
$\rho$ supersedes the background entirely at negative $x$ --- for every
$\rho>0$ --- and carries internal damage; $\rho=1$ is instead the
threshold at which the background ceases to reach the \emph{survivor's}
hazard, a statement about openness rather than about survival, since the
exponent is $\tfrac14$ for all $\rho>0$. The background,
built as a sum of exponentiated Brownian motions to force positivity, is
a partition function whose temperature falls with age: it freezes at
$2\log n/\epsilon^{2}$, is Gompertz-exponential before and
extreme-dominated after, and enters Section~\ref{sh:sec} only as a constant
headwind $\epsilon^{2}/2$ on the improvement rate --- gauge, hence
invisible. Strength and anatomy are not parameters but coordinates: the
state factors exactly and adaptedly as $x_t=\sigma t^{3/2}R_t\cos
\Theta_t$ with $R$ Rayleigh, $\Theta$ uniform and the two independent at
each fixed time, though \emph{not} as processes, the dependence peaking
at log-lag $\log3$; anatomy is phase within the excursion cycle, and the
system is in bad anatomy exactly half of logarithmic time. Between the
channels' shape walls lies an asymmetric waiting room $[-k/\rho,\,b]$
containing the origin --- which is the offset's purpose, since at $b=0$
the origin sits at the crest's inflection with gain $\tfrac12$ ---
within which the gain falls to $\kappa(\rho)e^{-(\rho b+k)/(1+\rho)}$
and the driver is accordingly muted, exponentially in the room's width;
the model leaves it in two ordered stages, internal at
$(k/\rho\sigma)^{2/3}$ and external at $(b/\sigma)^{2/3}$, of its own
accord, so no
capacity parameter is needed --- at the cost, recorded rather than
glossed, of revision~1's sharp threshold in an exponent, exchanged for a
crossover in a timescale plus a mechanical threshold at $\rho=1$. With
constant background the construction collapses to Section~\ref{sh:sec}'s model
with the credit passed through a soft kink, revealing the original
driver to have been the present driver times a state-dependent gain in
$(0,\rho)$ --- vanishing in the void, unit under shielding, $\rho$ under
damage. Both killing walls are $o(t^{3/2})$ and therefore gauge, so the
persistence exponent $\tfrac14$, the infinite life expectancy and the
survivorship interpolation survive exactly, with no truncation window:
the defect of the additive treatment is repaired, at the price of the
U-shaped hazard and its interior optimum, which no model satisfying
Constraint~\ref{tc:C3} can possess.
% =====================================================================
\PrintRefs

\end{document}
